% Generated by the blueprint tool (blueprint-pdf/1) from a blueprint snapshot; edit the
% Markdown sources, not this file.  Compile with `latexmk -pdf` or by running
% `pdflatex` twice (the second run settles the contents and the page references).
% Images, if any, are read from assets/ next to this file.
\documentclass[11pt]{article}
\usepackage[T1]{fontenc}
\usepackage[utf8]{inputenc}
\usepackage{lmodern}
\usepackage{textcomp}
\usepackage[a4paper,margin=2.5cm,headheight=14pt]{geometry}
\usepackage{amsmath,amssymb,amsthm}
\usepackage{amscd}  % KaTeX's CD diagrams; part of amsmath, so always there
\usepackage{mathtools}
\usepackage{bm}
\usepackage{graphicx}
\usepackage{xcolor}
\usepackage{array,booktabs,tabularx}
\usepackage{fancyvrb}
\usepackage[normalem]{ulem}
\usepackage{microtype}
% Optional, and only with their fonts: a small TeX installation has the
% packages without the fonts, which fails at the first formula.
\IfFileExists{mathrsfs.sty}{\IfFileExists{rsfs10.tfm}{\usepackage{mathrsfs}}{}}{}
\IfFileExists{stmaryrd.sty}{\IfFileExists{stmary10.tfm}{\usepackage{stmaryrd}}{}}{}
\IfFileExists{cancel.sty}{\usepackage{cancel}}{}
\usepackage[hyphens]{url}
\usepackage[pdfusetitle,bookmarksnumbered=false]{hyperref}
\usepackage{bookmark}
\definecolor{bplink}{HTML}{1F4E9C}
\definecolor{bpmuted}{HTML}{5F676E}
\hypersetup{colorlinks=true,linkcolor=bplink,urlcolor=bplink,citecolor=bplink}
\setlength{\parindent}{0pt}
\setlength{\parskip}{0.55em plus 0.2em minus 0.1em}
% Lists as tight as the page's.
\makeatletter
\def\@listi{\leftmargin\leftmargini\parsep 0pt\topsep 0.3em plus 0.1em\itemsep 0.1em plus 0.1em}
\let\@listI\@listi
\def\@listii{\leftmargin\leftmarginii\labelwidth\leftmarginii\advance\labelwidth-\labelsep\parsep 0pt\topsep 0.1em\itemsep 0.1em}
% An ordered list that does not start at 1.
\newcommand{\bpstartat}[1]{\setcounter{\@enumctr}{\numexpr#1-1\relax}}
\makeatother
\setcounter{tocdepth}{3}
\emergencystretch=3em
\allowdisplaybreaks

% A numbered statement, set as a paper sets one: "Theorem 2.1 (Title)."
% #1 is the lead ("Theorem~2.1"), #2 what follows it in roman.
\newenvironment{bpstatement}[2]{\par\addvspace{0.6\baselineskip}\noindent\textbf{#1}#2\textbf{.}\enspace\ignorespaces}
  {\par\addvspace{0.4\baselineskip}}
\newenvironment{bpstatementit}[2]{\begin{bpstatement}{#1}{#2}\itshape\ignorespaces}{\end{bpstatement}}
% A step: a relation with prose of its own, read under the object it starts from.
\newenvironment{bpstep}[1]{\par\addvspace{0.3\baselineskip}\begin{list}{}{\leftmargin=1.6em\rightmargin=0pt\topsep=0pt}\item\relax
  {\small\color{bpmuted}#1}\par\nopagebreak}{\end{list}\par}
% Notes about where an object is written out.
\newenvironment{bpnote}{\par\small\color{bpmuted}\itshape}{\par}
\newcommand{\bpomitted}[1]{\par{\small\color{bpmuted}\itshape #1\par}}
\newcommand{\bpbodyheading}[1]{\par\addvspace{0.4\baselineskip}\noindent\textbf{#1}\par\nopagebreak}
\newcommand{\bprule}{\par\medskip\noindent\hfil\rule{0.3\linewidth}{0.4pt}\hfil\par\medskip}
\newcommand{\bpboundary}[1]{\par\nopagebreak{\small\color{bpmuted}#1\par}\nopagebreak}
% A reference to something this export leaves out, and one that resolves nowhere.
\newcommand{\bpout}[1]{#1\textsuperscript{\hyperlink{bp-outside}{\ensuremath{\ast}}}}
\newcommand{\bpbroken}[1]{\textcolor{red!70!black}{[#1?]}}
\newcommand{\bpslug}[1]{\texttt{#1}}
% A bibliography entry: label, anchor, entry; the label hangs in the margin.
\newcommand{\bpbibitem}[3]{\par\setbox0\hbox{[#1]\enspace}\ifdim\wd0<5em\setbox0\hbox to 5em{[#1]\hfil}\fi
  \hangindent=\wd0 \hangafter=1 \noindent\box0\phantomsection\label{#2}#3\par\smallskip}
% A character the export knows no way to set.
\newcommand{\bpglyph}[1]{\fbox{\scriptsize #1}}
% Images, at their own size unless they are wider than the line.
\newsavebox{\bpimagebox}
\newcommand{\bpimage}[1]{\sbox{\bpimagebox}{\includegraphics{#1}}%
  \ifdim\wd\bpimagebox>\linewidth\resizebox{\linewidth}{!}{\usebox{\bpimagebox}}\else\usebox{\bpimagebox}\fi}
\newcommand{\bpmissingimage}[1]{\fbox{\parbox{0.8\linewidth}{\centering\small\color{bpmuted}Image not included: #1}}}
\DefineVerbatimEnvironment{bpverbatim}{Verbatim}{fontsize=\small,xleftmargin=1em}
\renewcommand{\qedsymbol}{\ensuremath{\blacksquare}}
\hypersetup{pdftitle={Weil Conjectures — 12 The trace formula},pdfsubject={Chapter 12 The trace formula; all levels},pdfcreator={blueprint}}
% KaTeX commands the website knows
\providecommand{\R}{\mathbb{R}}\providecommand{\N}{\mathbb{N}}\providecommand{\Z}{\mathbb{Z}}
\providecommand{\Q}{\mathbb{Q}}\providecommand{\C}{\mathbb{C}}
\providecommand{\Reals}{\mathbb{R}}\providecommand{\reals}{\mathbb{R}}\providecommand{\Complex}{\mathbb{C}}
\providecommand{\cnums}{\mathbb{C}}\providecommand{\natnums}{\mathbb{N}}\providecommand{\Bbb}[1]{\mathbb{#1}}
\providecommand{\bold}[1]{\mathbf{#1}}\providecommand{\frak}[1]{\mathfrak{#1}}\providecommand{\mathbbm}[1]{\mathbb{#1}}
\providecommand{\lang}{\langle}\providecommand{\rang}{\rangle}\providecommand{\bra}[1]{\langle #1\rvert}
\providecommand{\ket}[1]{\lvert #1\rangle}\providecommand{\braket}[1]{\langle #1\rangle}
\providecommand{\set}[1]{\{#1\}}\providecommand{\Set}[1]{\left\{#1\right\}}
\providecommand{\infin}{\infty}\providecommand{\isin}{\in}\providecommand{\plusmn}{\pm}\providecommand{\sub}{\subset}
\providecommand{\sube}{\subseteq}\providecommand{\supe}{\supseteq}\providecommand{\exist}{\exists}
\providecommand{\bull}{\bullet}\providecommand{\larr}{\leftarrow}\providecommand{\rarr}{\rightarrow}
\providecommand{\lrarr}{\leftrightarrow}\providecommand{\Larr}{\Leftarrow}\providecommand{\Rarr}{\Rightarrow}
\providecommand{\Lrarr}{\Leftrightarrow}\providecommand{\uarr}{\uparrow}\providecommand{\darr}{\downarrow}
\providecommand{\harr}{\leftrightarrow}\providecommand{\lArr}{\Leftarrow}\providecommand{\rArr}{\Rightarrow}
\providecommand{\hArr}{\Leftrightarrow}\providecommand{\alef}{\aleph}\providecommand{\alefsym}{\aleph}
\providecommand{\real}{\Re}\providecommand{\image}{\Im}\providecommand{\weierp}{\wp}
\providecommand{\Alpha}{\mathrm{A}}\providecommand{\Beta}{\mathrm{B}}\providecommand{\Epsilon}{\mathrm{E}}
\providecommand{\Zeta}{\mathrm{Z}}\providecommand{\Eta}{\mathrm{H}}\providecommand{\Iota}{\mathrm{I}}
\providecommand{\Kappa}{\mathrm{K}}\providecommand{\Mu}{\mathrm{M}}\providecommand{\Nu}{\mathrm{N}}
\providecommand{\Omicron}{\mathrm{O}}\providecommand{\Rho}{\mathrm{P}}\providecommand{\Tau}{\mathrm{T}}
\providecommand{\Chi}{\mathrm{X}}\providecommand{\omicron}{o}\providecommand{\gt}{>}\providecommand{\lt}{<}
\providecommand{\And}{\&}\providecommand{\hdashline}{\hline}\providecommand{\operatornamewithlimits}[1]{\operatorname*{#1}}
\providecommand{\llbracket}{[\![}\providecommand{\rrbracket}{]\!]}\providecommand{\cancel}[1]{#1}
\providecommand{\mathscr}{\mathcal}
% the characters beyond ASCII this document uses
\DeclareUnicodeCharacter{2160}{\textup{I}}
\DeclareUnicodeCharacter{2161}{\textup{II}}
\DeclareUnicodeCharacter{2162}{\textup{III}}
% the project's macros ([katex.macros] in blueprint.toml)
\providecommand{\A}{}\renewcommand{\A}{\mathbf A}
\providecommand{\AEn}{}\renewcommand{\AEn}{A_{E,n}}
\providecommand{\Aaf}{}\renewcommand{\Aaf}{\mathbf A}
\providecommand{\Ab}{}\renewcommand{\Ab}{\mathsf{Ab}}
\providecommand{\AffSch}{}\renewcommand{\AffSch}{\operatorname{AffSch}}
\providecommand{\Aone}{}\renewcommand{\Aone}{\mathbf A^1}
\providecommand{\Bl}{}\renewcommand{\Bl}{\operatorname{Bl}}
\providecommand{\Br}{}\renewcommand{\Br}{\mathrm{Br}}
\providecommand{\C}{}\renewcommand{\C}{\mathbb C}
\providecommand{\Cech}{}\renewcommand{\Cech}{\check{H}}
\providecommand{\Coext}{}\renewcommand{\Coext}{\operatorname{Coext}}
\providecommand{\Cone}{}\renewcommand{\Cone}{\operatorname{Cone}}
\providecommand{\Cons}{}\renewcommand{\Cons}{\operatorname{Cons}}
\providecommand{\Cref}{}\renewcommand{\Cref}[1]{\text{#1}}
\providecommand{\D}{}\renewcommand{\D}{D}
\providecommand{\DLAFT}{}\renewcommand{\DLAFT}{\operatorname{FT}}
\providecommand{\DLAFr}{}\renewcommand{\DLAFr}{\operatorname{Fr}}
\providecommand{\DLAQ}{}\renewcommand{\DLAQ}{{\overline{\mathbf Q}_\ell}}
\providecommand{\DLRQ}{}\renewcommand{\DLRQ}{{\overline{\mathbf Q}_\ell}}
\providecommand{\Dbc}{}\renewcommand{\Dbc}{D^b_c}
\providecommand{\Dc}{}\renewcommand{\Dc}{D_{\mathrm{cons}}}
\providecommand{\Dctf}{}\renewcommand{\Dctf}{D^b_{ctf}}
\providecommand{\Dic}{}\renewcommand{\Dic}{D_{\IndCons}}
\providecommand{\Dplus}{}\renewcommand{\Dplus}{D^{+}}
\providecommand{\End}{}\renewcommand{\End}{\operatorname{End}}
\providecommand{\Et}{}\renewcommand{\Et}{\mathrm{\acute{e}t}}
\providecommand{\Ex}{}\renewcommand{\Ex}{\mathrm{Ex}}
\providecommand{\Ext}{}\renewcommand{\Ext}{\operatorname{Ext}}
\providecommand{\F}{}\renewcommand{\F}{\mathbb F}
\providecommand{\FT}{}\renewcommand{\FT}{\operatorname{FT}}
\providecommand{\Fbar}{}\renewcommand{\Fbar}{{\overline{\mathbf F}_q}}
\providecommand{\Fell}{}\renewcommand{\Fell}{\mathbf F_\ell}
\providecommand{\Fib}{}\renewcommand{\Fib}{\operatorname{Fib}}
\providecommand{\Fix}{}\renewcommand{\Fix}{\operatorname{Fix}}
\providecommand{\Fp}{}\renewcommand{\Fp}{{\mathbf F_{p}}}
\providecommand{\Fq}{}\renewcommand{\Fq}{{\mathbf F_q}}
\providecommand{\Fqn}{}\renewcommand{\Fqn}{{\mathbf F_{q^{n}}}}
\providecommand{\Frob}{}\renewcommand{\Frob}{\operatorname{Frob}}
\providecommand{\GL}{}\renewcommand{\GL}{\operatorname{GL}}
\providecommand{\Gal}{}\renewcommand{\Gal}{\operatorname{Gal}}
\providecommand{\Hens}{}\renewcommand{\Hens}{\operatorname{Hens}}
\providecommand{\Hilb}{}\renewcommand{\Hilb}{\operatorname{Hilb}}
\providecommand{\Hom}{}\renewcommand{\Hom}{\operatorname{Hom}}
\providecommand{\IndCons}{}\renewcommand{\IndCons}{\operatorname{IndCons}}
\providecommand{\Kar}{}\renewcommand{\Kar}{\operatorname{Kar}}
\providecommand{\LT}{}\renewcommand{\LT}{\operatorname{LT}}
\providecommand{\Lie}{}\renewcommand{\Lie}{\operatorname{Lie}}
\providecommand{\Loc}{}\renewcommand{\Loc}{\operatorname{Loc}}
\providecommand{\Mod}{}\renewcommand{\Mod}{\mathsf{Mod}}
\providecommand{\N}{}\renewcommand{\N}{\mathbb{N}}
\providecommand{\Needspace}{}\renewcommand{\Needspace}{\par}
\providecommand{\O}{}\renewcommand{\O}{{\mathcal{O}}}
\providecommand{\Ob}{}\renewcommand{\Ob}{\mathrm{Ob}}
\providecommand{\Oe}{}\renewcommand{\Oe}{\mathcal O_E}
\providecommand{\PP}{}\renewcommand{\PP}{\mathbf{P}}
\providecommand{\Perf}{}\renewcommand{\Perf}{\operatorname{Perf}}
\providecommand{\Pic}{}\renewcommand{\Pic}{\mathrm{Pic}}
\providecommand{\Pm}{}\renewcommand{\Pm}{\mathbf P^{m}}
\providecommand{\Pone}{}\renewcommand{\Pone}{\mathbf P^{1}}
\providecommand{\Pqn}{}\renewcommand{\Pqn}{\mathbf P^{n}}
\providecommand{\ProFin}{}\renewcommand{\ProFin}{\operatorname{ProFin}}
\providecommand{\Proet}{}\renewcommand{\Proet}{\mathrm{pro\acute{e}t}}
\providecommand{\Q}{}\renewcommand{\Q}{\mathbb Q}
\providecommand{\Qb}{}\renewcommand{\Qb}{\overline{\mathbf Q}_{\ell}}
\providecommand{\Qbar}{}\renewcommand{\Qbar}{\overline{\mathbf Q}}
\providecommand{\QbarX}{}\renewcommand{\QbarX}{\overline{\mathbf Q}_{\ell,X}}
\providecommand{\Qbarl}{}\renewcommand{\Qbarl}{{\overline{\mathbf Q}_{\ell}}}
\providecommand{\Qell}{}\renewcommand{\Qell}{{\mathbf{Q}_{\ell}}}
\providecommand{\Ql}{}\renewcommand{\Ql}{\mathbf Q_{\ell}}
\providecommand{\Qlb}{}\renewcommand{\Qlb}{{\overline{\mathbf Q}_\ell}}
\providecommand{\Qlbar}{}\renewcommand{\Qlbar}{\overline{\mathbf Q}_\ell}
\providecommand{\R}{}\renewcommand{\R}{\mathbb{R}}
\providecommand{\RG}{}\renewcommand{\RG}{\operatorname{\mathrm{R}\Gamma}}
\providecommand{\RGamma}{}\renewcommand{\RGamma}{R\Gamma}
\providecommand{\RGammac}{}\renewcommand{\RGammac}{R\Gamma_c}
\providecommand{\RHom}{}\renewcommand{\RHom}{R\mathcal{H}om}
\providecommand{\RHomX}{}\renewcommand{\RHomX}{R\mathcal H om_X}
\providecommand{\Rat}{}\renewcommand{\Rat}{\operatorname{IsRat}}
\providecommand{\Res}{}\renewcommand{\Res}{\operatorname{Res}}
\providecommand{\Rlim}{}\renewcommand{\Rlim}{R\!\varprojlim}
\providecommand{\SIm}{}\renewcommand{\SIm}{\operatorname{SIm}}
\providecommand{\Sch}{}\renewcommand{\Sch}{\mathsf{Sch}}
\providecommand{\Sh}{}\renewcommand{\Sh}{\operatorname{Sh}}
\providecommand{\Span}{}\renewcommand{\Span}{\operatorname{span}}
\providecommand{\Spec}{}\renewcommand{\Spec}{\operatorname{Spec}}
\providecommand{\Tr}{}\renewcommand{\Tr}{\operatorname{Tr}}
\providecommand{\Trsm}{}\renewcommand{\Trsm}{\operatorname{Tr}^{\mathrm{sm}}}
\providecommand{\Wt}{}\renewcommand{\Wt}{\operatorname{Wt}}
\providecommand{\Z}{}\renewcommand{\Z}{\mathbb{Z}}
\providecommand{\Zar}{}\renewcommand{\Zar}{\mathrm{zar}}
\providecommand{\Zell}{}\renewcommand{\Zell}{{\mathbf{Z}_{\ell}}}
\providecommand{\Zl}{}\renewcommand{\Zl}{{\mathbf Z_\ell}}
\providecommand{\affet}{}\renewcommand{\affet}[1]{{#1}^{\mathrm{aff}}_{\mathrm{ét}}}
\providecommand{\affproet}{}\renewcommand{\affproet}[1]{{#1}^{\mathrm{aff}}_{\mathrm{proét}}}
\providecommand{\boxt}{}\renewcommand{\boxt}{\boxtimes_A^{\mathrm L}}
\providecommand{\cF}{}\renewcommand{\cF}{\mathcal{F}}
\providecommand{\cHom}{}\renewcommand{\cHom}{R\mathcal{H}\!om}
\providecommand{\cI}{}\renewcommand{\cI}{\mathcal{I}}
\providecommand{\cL}{}\renewcommand{\cL}{\mathcal L}
\providecommand{\cM}{}\renewcommand{\cM}{\mathcal M}
\providecommand{\cN}{}\renewcommand{\cN}{\mathcal N}
\providecommand{\cO}{}\renewcommand{\cO}{\mathcal O}
\providecommand{\cT}{}\renewcommand{\cT}{\mathcal T}
\providecommand{\checked}{}\renewcommand{\checked}{\color{red}\aftergroup\restorecolor}
\providecommand{\coeff}{}\renewcommand{\coeff}{\operatorname{coeff}}
\providecommand{\coker}{}\renewcommand{\coker}{\operatorname{coker}}
\providecommand{\colim}{}\renewcommand{\colim}{\operatorname{colim}}
\providecommand{\constructedby}{}\renewcommand{\constructedby}[1]{\par\smallskip\noindent\textcolor{purple!60!black}{\textsf{Constructed by:} }
  \footnotesize\sloppy\nolinkurl{#1}\par}
\providecommand{\deepresult}{}\renewcommand{\deepresult}[1]{\par\smallskip\noindent\textcolor{orange!45!black}{
  \textsf{Blueprint role: named deep theorem node; see the cluster proof architecture.}}
  \par\noindent\textcolor{orange!45!black}{\textsf{Mathematical source:} }#1\par}
\providecommand{\dt}{}\renewcommand{\dt}{\mathbin{\otimes^{\mathbf L}}}
\providecommand{\dw}{}\renewcommand{\dw}{\operatorname{dw}}
\providecommand{\et}{}\renewcommand{\et}[1]{{#1}_{\mathrm{ét}}}
\providecommand{\ev}{}\renewcommand{\ev}{\operatorname{ev}}
\providecommand{\fib}{}\renewcommand{\fib}{\operatorname{fib}}
\providecommand{\formalizationnote}{}\renewcommand{\formalizationnote}[1]{\begin{quote}\small\raggedright\textbf{Formalization note.} #1\end{quote}}
\providecommand{\fpqc}{}\renewcommand{\fpqc}{X_{\mathrm{fpqc}}}
\providecommand{\gammacont}{}\renewcommand{\gammacont}{\Gamma_{\mathrm{cont}}}
\providecommand{\gm}{}\renewcommand{\gm}{\mathbf G_m}
\providecommand{\hocolim}{}\renewcommand{\hocolim}{\operatorname{hocolim}}
\providecommand{\holim}{}\renewcommand{\holim}{\operatorname{holim}}
\providecommand{\hrB}{}\renewcommand{\hrB}{\@ifnextchar]{\hrEnd}{\hrStep}}
\providecommand{\hrEnd}{}\renewcommand{\hrEnd}{\@firstoftwo{}}
\providecommand{\hrStep}{}\renewcommand{\hrStep}{\@firstoftwo{\hrB}}
\providecommand{\id}{}\renewcommand{\id}{\operatorname{id}}
\providecommand{\kunFl}{}\renewcommand{\kunFl}{\mathbf F_\ell}
\providecommand{\kunZn}{}\renewcommand{\kunZn}{\mathbf Z/n\mathbf Z}
\providecommand{\kuncd}{}\renewcommand{\kuncd}{\operatorname{cd}}
\providecommand{\kundim}{}\renewcommand{\kundim}{\dim}
\providecommand{\kunotimes}{}\renewcommand{\kunotimes}[1]{\otimes^{\mathbf L}_{#1}}
\providecommand{\kunsupp}{}\renewcommand{\kunsupp}{\operatorname{Supp}}
\providecommand{\lotimes}{}\renewcommand{\lotimes}{\otimes^{\mathrm L}}
\providecommand{\m}{}\renewcommand{\m}{{\mathfrak{m}}}
\providecommand{\n}{}\renewcommand{\n}{{\mathfrak{n}}}
\providecommand{\nagataGap}{}\renewcommand{\nagataGap}[1]{\par\noindent\textbf{Formalization gap.} #1\par}
\providecommand{\openobligation}{}\renewcommand{\openobligation}[1]{\par\smallskip\noindent\textcolor{red!65!black}{\textsf{Open obligation:} }
  \footnotesize\sloppy\nolinkurl{#1}\par}
\providecommand{\p}{}\renewcommand{\p}{{\mathfrak{p}}}
\providecommand{\pr}{}\renewcommand{\pr}{\operatorname{pr}}
\providecommand{\preshv}{}\renewcommand{\preshv}[1]{\mathrm{PreShv}(#1)}
\providecommand{\pro}{}\renewcommand{\pro}[1]{\mathrm{Pro}(#1)}
\providecommand{\proet}{}\renewcommand{\proet}{X_{\mathrm{proét}}}
\providecommand{\q}{}\renewcommand{\q}{{\mathfrak{q}}}
\providecommand{\rad}{}\renewcommand{\rad}{\operatorname{rad}}
\providecommand{\rank}{}\renewcommand{\rank}{\operatorname{rank}}
\providecommand{\red}{}\renewcommand{\red}{\color{red}}
\providecommand{\restorecolor}{}\renewcommand{\restorecolor}{\color{black}}
\providecommand{\retiredobligation}{}\renewcommand{\retiredobligation}[1]{\par\smallskip\noindent\textcolor{red!75!black}{\textsf{Retired invalid route:} }
  \footnotesize\sloppy\nolinkurl{#1}\par}
\providecommand{\rk}{}\renewcommand{\rk}{\operatorname{rk}}
\providecommand{\sep}{}\renewcommand{\sep}{\textnormal{sep}}
\providecommand{\shv}{}\renewcommand{\shv}[1]{\mathrm{Shv}(#1)}
\providecommand{\spec}{}\renewcommand{\spec}[1]{\mathrm{Spec}(#1)}
\providecommand{\stacks}{}\renewcommand{\stacks}[1]{\href{https://stacks.math.columbia.edu/tag/#1}{\textnormal{Stacks, Tag~#1}}}
\providecommand{\stacksproject}{}\renewcommand{\stacksproject}[1]{{\cite[\href{https://stacks.math.columbia.edu/tag/#1}{Tag #1}]{stacks-project}}}
\providecommand{\statusnote}{}\renewcommand{\statusnote}[1]{\begin{quote}\small\raggedright\textbf{Status.} #1\end{quote}}
\providecommand{\tildeideal}{}\renewcommand{\tildeideal}[2]{\widetilde{(#1, #2)}}
\providecommand{\tildering}{}\renewcommand{\tildering}[3]{{#1}_{\widetilde{(#2, #3)}}}
\providecommand{\tr}{}\renewcommand{\tr}{\operatorname{tr}}
\providecommand{\wt}{}\renewcommand{\wt}{\operatorname{wt}}

\begin{document}
\pagestyle{plain}
\begin{center}
{\LARGE\bfseries Weil Conjectures\par}
\vspace{0.7em}
{\large Chapter 12: The trace formula\par}
\vspace{0.5em}
{\small\color{bpmuted}Snapshot 2026-09-26, revision 8e4cb88e2a\par}
\end{center}
\vspace{0.5em}
{\small\tableofcontents}
\bigskip\hrule\bigskip

% object: chap-The-trace-formula (blueprint/chap-The-trace-formula/_section.md)
\section*{\texorpdfstring{12\quad The trace formula}{12 The trace formula}}
\phantomsection\label{bp1}%
\addcontentsline{toc}{section}{\texorpdfstring{12\quad The trace formula}{12 The trace formula}}
% text: prose-12-01
\phantomsection\label{bp2}%
Let \(X_0\) be a separated scheme of finite type over \(\mathbf F_q\), let \(X=X_0\times_{\mathbf F_q}\overline{\mathbf F}_q\), and let \(\ell\nmid q\). By Lemma \hyperref[bp15]{12.5}, the points of \(X_0\) over \(\mathbf F_{q^m}\) are the fixed points of the \(m\)-th power of the Frobenius \(\Phi_X\) on \(X\). The Lefschetz–Verdier formula (Theorem \bpout{11.13}) writes a trace on cohomology as a sum of local terms, but it says nothing about what a local term is. To count points we need the local term of Frobenius at a fixed point to be the trace on the stalk.

For a self-map \(\varphi\) of a compact manifold with transverse fixed points, the Lefschetz formula counts each fixed point \(x\) with the sign of \(\det(1-d\varphi_x)\). The differential of Frobenius is zero, so we expect every sign to be \(+1\). On a curve this can be seen directly. If \(y\) is a fixed point of \(\Phi_X^m\), \(u\) is a uniformizer at \(y\) defined over \(\mathbf F_{q^m}\), and \(v\) is the same function on the second factor, then near \((y,y)\) the graph of \(\Phi_X^m\) is \(v=u^{q^m}\), and it meets the diagonal \(v=u\) with multiplicity one, since \(u^{q^m}-u\) has a simple zero at \(u=0\).

% object: question-lefschetz-formula-frobenius (blueprint/chap-The-trace-formula/question-lefschetz-formula-frobenius/_section.md)
\phantomsection\label{bp3}%
\begin{bpstatement}{Question~12.1}{ (The Lefschetz formula for Frobenius)}
Does the Lefschetz formula hold for Frobenius?
\end{bpstatement}

% text: prose-12-02
\phantomsection\label{bp4}%
For constant coefficients the local term at a transverse fixed point is \(1\), so on a proper \(X\) transversality would be enough. Two things stand in the way in general. First, \(X\) need not be proper. If we compactify, the new fixed points at the boundary have zero stalk for \(j_!K\), and we need their local terms to vanish. Second, a local term is computed from \(\cHom(c_2^*K,c_1^!K)\) near the fixed point, where \(c_1,c_2\) are the two projections of the graph of Frobenius to \(X\) (Definition \bpout{11.5}). The stalk \(K_{\bar y}\) does not determine this complex (Remark \hyperref[bp14]{12.4}), so the vanishing of the local term at a point where the stalk vanishes is a real statement, which transversality alone does not give.

One computation on a curve settles both problems: Proposition \hyperref[bp49]{12.16} says that for any morphism \(\psi:(\Phi_X^m)^*K\to K\) on a smooth curve, the local term at a rational point \(y\) is \(\operatorname{tr}(\psi_{\bar y}\mid K_{\bar y})\). The proof, which is Grothendieck's [\hyperref[bp105]{SGA5}, XV, §3], uses \(d\Phi_X=0\) a second time: the graph \(v=u^{q^m}\) is tangent to the axis \(v=0\) and the diagonal is not, so blowing up \((y,y)\) separates them, and near the exceptional divisor the local term becomes a restriction map that vanishes when \(K_{\bar y}=0\) (Remark \hyperref[bp18]{12.6}). Since both sides of the formula are additive along an open subscheme and its complement (Lemma \hyperref[bp58]{12.19}) and compatible with \(Rf_!\) (Lemma \hyperref[bp55]{12.18}), we never need to compactify anything but \(\mathbf A^1\subset\mathbf P^1\).

The answer is Theorem \hyperref[bp73]{12.24}: for every Weil complex \(K\in D^b_{\mathrm{ctf}}(X,\Lambda)\) (Section \bpout{10.5}), where \(\Lambda\in\{A,\mathbf Z_\ell,\mathcal O_E,\mathbf Q_\ell,E,\overline{\mathbf Q}_\ell\}\) and \(A\) is a finite commutative \(\ell\)-primary ring, and for every \(m\geq1\), we have in \(\Lambda\)

\[
\operatorname{tr}_\Lambda\bigl(\operatorname{Frob}_{X,K}^m\mid
R\Gamma_c(X_{\mathrm{pro}\acute{e}t},K)\bigr)
=\sum_{y\in X_0(\mathbf F_{q^m})}
\operatorname{tr}_\Lambda\bigl(\operatorname{Frob}_{y,m}\mid
K_{\bar y}\bigr).
\]

Here \(\operatorname{tr}_\Lambda\) is the trace of an endomorphism of a perfect complex (Definition \bpout{11.1}), and \(\operatorname{Frob}_{y,m}\) is the \(m\)-th iterate of the Weil structure on the stalk at the fixed point \(\bar y\). For \(K=\Lambda_X\) and \(\Lambda\) a field the formula is Weil's axiom (b), and it gives the rationality of the zeta function (Corollary \hyperref[bp95]{12.31}).

Purity, however, needs the formula for general \(K\). For \(\Lambda\in\{\mathbf Q_\ell,E,\overline{\mathbf Q}_\ell\}\), in the form of an Euler product (Theorem \hyperref[bp92]{12.30}), it reads

\[
\prod_{x\in|X_0|}\det\bigl(1-T^{\deg x}\operatorname{Frob}_x\mid K_{\bar x}\bigr)^{-1}
=\prod_i\det\bigl(1-T\operatorname{Frob}_{X,K}\mid H^i_c(X,K)\bigr)^{(-1)^{i+1}},
\]

which relates the local eigenvalues of \(K\) to the eigenvalues on cohomology.

The conventions of Chapter \bpout{11} remain in force.

% object: sec-Frobenius-correspondences (blueprint/chap-The-trace-formula/sec-Frobenius-correspondences/_section.md), content hidden
\subsection*{\texorpdfstring{12.1\quad Frobenius correspondences}{12.1 Frobenius correspondences}}
\phantomsection\label{bp5}%
\addcontentsline{toc}{subsection}{\texorpdfstring{12.1\quad Frobenius correspondences}{12.1 Frobenius correspondences}}
% text: prose-12-03
\phantomsection\label{bp6}%
Two Frobenius maps act on \(X=X_0\times_{k_0}k\): the \(k\)-linear scheme Frobenius \(\Phi_X\) and the semilinear geometric Frobenius \(\operatorname{Fr}_X\) that defines Weil structures. We prove the trace formula for \(\Phi_X\), whose graph is a correspondence over \(k\), and Lemma \hyperref[bp9]{12.3} translates it to \(\operatorname{Fr}_X\).

The operator \(T_{c_m}\) on cohomology with compact support defined below makes sense even when \(X\) is not proper, because the second leg \(\Phi_X^m\) of the graph \(c_m\) is proper. Indeed, \(\Phi_X\) and all its iterates are finite, since a finitely generated algebra over a perfect field is finite over its subring of \(q\)-th powers.

We cannot work with \(\operatorname{Fr}_X\) directly. Indeed, the Weil structure \(\phi_K\) is a map \(\operatorname{Fr}_X^*K\to K\), but \(\operatorname{Fr}_X\) lies over the automorphism \(\operatorname{Spec}(\sigma)\) of \(s\). Hence its graph is not a morphism \(X\to X\times_sX\) of \(k\)-schemes. Since the formalism of Chapter \bpout{11} uses \(c_1^!\) only for \(k\)-morphisms, it does not apply to this graph. The \(k\)-linear \(\Phi_X=F_X\circ\operatorname{Fr}_X\), on the other hand, has a graph of the required kind. Pullback along the absolute Frobenius \(F_X\) is the identity on étale sheaves, and this is how Lemma \hyperref[bp9]{12.3} turns \(\phi_K\) into a map \(\Phi_X^*K\to K\). For example, on \(X=\mathbf A^1_k\) the map \(\Phi_X\) is \(t\mapsto t^q\), and the fixed points of \(\Phi_X^m\) are the \(q^m\) elements of \(\mathbf F_{q^m}\). Lemma \hyperref[bp15]{12.5} describes them in general.

We allow maps \(\Phi_X^*K\to K\) that are not isomorphisms, because the proof over \(E\) clears denominators: it multiplies a Weil structure by a power of \(\pi_E\) (Lemma \hyperref[bp69]{12.22}). The operators below are the two sides of the trace formula in Theorem \hyperref[bp62]{12.20}.

% object: definition-frobenius-correspondences-their-operators (blueprint/chap-The-trace-formula/sec-Frobenius-correspondences/definition-frobenius-correspondences-their-operators/_section.md)
\phantomsection\label{bp7}%
\begin{bpstatement}{Definition~12.2}{ (Frobenius correspondences and their operators)}
Let \(X_0/k_0\) be separated and of finite type, and put \(X=X_0\times_{k_0}k\). Denote by \(\Phi_X:X\to X\) the base change of the \(q\)-power scheme Frobenius of \(X_0\). It is a finite \(k\)-morphism and a universal homeomorphism. A \emph{Frobenius correspondence} on \(K\) is a morphism, not necessarily invertible,

\[
\psi:\Phi_X^*K\longrightarrow K,
\qquad K\in D^b_{\mathrm{ctf}}(X,A).
\]

Its iterates are the morphisms

\[
\psi_m=\psi\circ\Phi_X^*\psi\circ\cdots\circ
(\Phi_X^{m-1})^*\psi:(\Phi_X^m)^*K\longrightarrow K.
\]

Let \(c_m=(\operatorname{id}_X,\Phi_X^m):X\to X\times_sX\) be the graph of \(\Phi_X^m\). Since its first leg is the identity, \(\psi_m\) is a cohomological correspondence on \((X,K)\) for \(c_m\). Its operator on cohomology with compact support (Definition \bpout{11.11}) is the endomorphism

\[
T_{c_m}(\psi_m)
=R\Gamma_c(X,\psi_m)\circ(\Phi_X^m)^*,
\]

where \((\Phi_X^m)^*:R\Gamma_c(X,K)\to R\Gamma_c(X,(\Phi_X^m)^*K)\) is the pullback map induced by the unit of \((\Phi_X^m)^*\dashv R(\Phi_X^m)_*\). For \(x\in X_0(\mathbf F_{q^m})\) let \(\bar x:s\to X\) be the associated fixed point. The stalk of \(\psi_m\) at \(\bar x\) is an endomorphism \(\psi_{m,\bar x}:K_{\bar x}\to K_{\bar x}\).
\end{bpstatement}

% text: prose-12-04
\phantomsection\label{bp8}%
The operator \(T_{c_m}(\psi_m)\) is defined with the \(k\)-linear \(\Phi_X\), while the Frobenius operator of Section \bpout{10.5}, whose trace the theorem is about, is defined with \(\operatorname{Fr}_X\). The next lemma says that the two agree, so that the trace formula proved for \(\Phi_X\) is a statement about \(\operatorname{Frob}_{X,K}\).

% object: lemma-scheme-frobenius-geometric-frobenius (blueprint/chap-The-trace-formula/sec-Frobenius-correspondences/lemma-scheme-frobenius-geometric-frobenius/_section.md)
\phantomsection\label{bp9}%
\begin{bpstatementit}{Lemma~12.3}{ (Scheme Frobenius and geometric Frobenius)}
Let \(\operatorname{Fr}_X\) be the semilinear geometric Frobenius of Definition \bpout{10.18}, and for a morphism \(\phi:\operatorname{Fr}_X^*K\to K\) put \(\phi_m=\phi\circ \operatorname{Fr}_X^*\phi\circ\cdots\circ(\operatorname{Fr}_X^{m-1})^*\phi\). There are natural isomorphisms

\[
\kappa_m:(\Phi_X^m)^*K\xrightarrow{\sim}(\operatorname{Fr}_X^m)^*K
\]

compatible with iteration. For a Weil structure \(\phi:\operatorname{Fr}_X^*K\xrightarrow{\sim}K\), put \(\psi=\phi\kappa_1\). Then

\[
\psi_m=\phi_m\kappa_m,
\qquad T_{c_m}(\psi_m)=\operatorname{Frob}_{X,K}^{m},
\]

where \(\operatorname{Frob}_{X,K}\) is the geometric Frobenius operator on compact cohomology of Definition \bpout{10.19}. The same identity of operators holds for an arbitrary morphism \((\operatorname{Fr}_X^m)^*K\to K\) in place of \(\phi_m\).
\end{bpstatementit}

% proof: proof/lemma-scheme-frobenius-geometric-frobenius
\begin{proof}[{Proof}]\phantomsection\label{bp10}%
The maps \(\Phi_X\) and \(\operatorname{Fr}_X\) differ by the absolute Frobenius, and pullback along the absolute Frobenius is naturally the identity on étale sheaves. This gives \(\kappa_m\), and we then compare the two pullback maps on one compactification. On coordinate rings the two maps are

\[
\Phi_X^\#(a\otimes b)=a^q\otimes b,
\qquad \operatorname{Fr}_X^\#(a\otimes b)=a\otimes b^{1/q}
\]

so that \(\Phi_X=F_X\operatorname{Fr}_X\), where \(F_X\) is the absolute \(q\)-power Frobenius of \(X\). The absolute Frobenius commutes with every morphism of schemes over \(\mathbf F_p\), so also \(\Phi_X^m=\operatorname{Fr}_X^mF_X^m\). On the classical constructible category, pullback along absolute Frobenius is naturally isomorphic to the identity, because for an étale object its relative Frobenius is an isomorphism, functorial in that object. Transporting this natural isomorphism to \(D^b_{\mathrm{ctf}}(X,A)\) by the comparison with the pro-étale site, we obtain \(\theta:F_X^*\simeq\mathrm{id}\). It is compatible with pullback along every morphism \(h\), that is, \(h^*\theta=\theta h^*\) under \(h^*F^*=F^*h^*\), since relative Frobenius of étale objects is compatible with base change. Let \(\theta^m:(F_X^m)^*\simeq\mathrm{id}\) be its \(m\)-fold composite and put

\[
\kappa_m=\theta^m_{(\operatorname{Fr}_X^m)^*K}:
(\Phi_X^m)^*K=(F_X^m)^*(\operatorname{Fr}_X^m)^*K
\xrightarrow{\sim}(\operatorname{Fr}_X^m)^*K .
\]

Compatibility with iteration means that \(\kappa_{m+1}=\kappa_{m,\operatorname{Fr}_X^*K}\circ(\Phi_X^m)^*\kappa_1\), and it follows from the compatibility of \(\theta\) with pullback along \(\Phi_X^m\). Now \(\psi_m=\phi_m\kappa_m\) follows by induction on \(m\). If it holds for \(m\), then

\[
\begin{aligned}
\psi_{m+1}&=\psi_m\circ(\Phi_X^m)^*\psi
=\phi_m\circ\kappa_m\circ(\Phi_X^m)^*\phi\circ(\Phi_X^m)^*\kappa_1\\
&=\phi_m\circ(\operatorname{Fr}_X^m)^*\phi\circ\kappa_{m,\operatorname{Fr}_X^*K}\circ(\Phi_X^m)^*\kappa_1
=\phi_{m+1}\kappa_{m+1},
\end{aligned}
\]

where the third equality is the naturality of \(\kappa_m\) in \(K\).

For the pullback maps on compact cohomology, let \(j:X\hookrightarrow\bar X\) be the base change of a compactification of \(X_0\) over \(k_0\), and use the pullback maps \(\operatorname{pb}_h\) of Section \bpout{10.5}. The endomorphisms \(\Phi_{\bar X}\), \(\operatorname{Fr}_{\bar X}\) and \(F_{\bar X}\) of the scheme \(\bar X\) restrict to \(\Phi_X\), \(\operatorname{Fr}_X\) and \(F_X\). Each of them has inverse image \(X\) of \(X\), so each commutes with \(j_!\), and \(\Phi_{\bar X}^m=\operatorname{Fr}_{\bar X}^mF_{\bar X}^m\). Put \(L=j_!K\), so that \(R\Gamma_c(X,K)=R\Gamma(\bar X,L)\).

First, the map \((\Phi_X^m)^*\) of Definition \hyperref[bp7]{12.2} is \(\operatorname{pb}_{\Phi_{\bar X}^m,L}\). Indeed, \(\Phi_X^m\) is finite, so \(R(\Phi_X^m)_!=R(\Phi_X^m)_*\), and the square \(j\Phi_X^m=\Phi_{\bar X}^mj\) is cartesian, so \(j_!R(\Phi_X^m)_*\simeq R(\Phi_{\bar X}^m)_*j_!\). Under this isomorphism the identification \(R\Gamma_c(X,R(\Phi_X^m)_*-)=R\Gamma_c(X,-)\) used there is the composition of direct images \(R\Gamma(\bar X,R(\Phi_{\bar X}^m)_*-)=R\Gamma(\bar X,-)\), and the unit for \(\Phi_X^m\) corresponds to the unit for \(\Phi_{\bar X}^m\).

Second, by transitivity, \(\operatorname{pb}_{\Phi_{\bar X}^m,L}
 =\operatorname{pb}_{F_{\bar X}^m,(\operatorname{Fr}_{\bar X}^m)^*L}\circ
 \operatorname{pb}_{\operatorname{Fr}_{\bar X}^m,L}\). The absolute Frobenius \(F_{\bar X}^m\) satisfies the hypothesis of [Stacks, \href{https://stacks.math.columbia.edu/tag/03SP}{03SP}], since its relative Frobenius over each étale object is an isomorphism, functorial in that object. Since \(A\) is finite, we may compute on the étale site (Proposition \bpout{5.15}). There the formal argument of [Stacks, \href{https://stacks.math.columbia.edu/tag/03SP}{03SP}] identifies \((F_{\bar X}^m)_*\) and \((F_{\bar X}^m)^{-1}\) with the identity functors of the topos, compatibly with the unit, and this identification is \(\theta^m\), since both come from the relative Frobenius isomorphisms of étale objects. Hence for every \(N\in D^b_{\mathrm{ctf}}(\bar X,A)\) the pullback map \(\operatorname{pb}_{F_{\bar X}^m,N}\) followed by \(R\Gamma(\bar X,\theta^m_N)\) is the identity of \(R\Gamma(\bar X,N)\). With \(N=(\operatorname{Fr}_{\bar X}^m)^*L\) this gives

\[
R\Gamma_c(X,\kappa_m)\circ(\Phi_X^m)^*=\operatorname{pb}_{\operatorname{Fr}_X^m} .
\]

Finally, the naturality of \(\operatorname{pb}_{\operatorname{Fr}_X}\) moves the maps \(R\Gamma_c(X,(\operatorname{Fr}_X^i)^*\phi)\) past the pullback maps, and transitivity gives \(\operatorname{Frob}_{X,K}^m=R\Gamma_c(X,\phi_m)\circ\operatorname{pb}_{\operatorname{Fr}_X^m}\). Hence

\[
T_{c_m}(\psi_m)=R\Gamma_c(X,\phi_m)\circ R\Gamma_c(X,\kappa_m)\circ(\Phi_X^m)^*
=R\Gamma_c(X,\phi_m)\circ\operatorname{pb}_{\operatorname{Fr}_X^m}
=\operatorname{Frob}_{X,K}^m .
\]

The same computation with an arbitrary morphism \(\chi:(\operatorname{Fr}_X^m)^*K\to K\) in place of \(\phi_m\) gives \(T_{c_m}(\chi\kappa_m)=R\Gamma_c(X,\chi)\circ\operatorname{pb}_{\operatorname{Fr}_X^m}\).
\end{proof}

% text: prose-12-05
\phantomsection\label{bp11}%
For example, if \(X_0\) is smooth and geometrically connected of dimension \(d\) and \(K=A_X\) with its constant Weil structure, then \(T_{c_1}(\psi)\) acts on \(H^{2d}_c(X,A)=A(-d)\) by \(q^d\), which agrees with the action of geometric Frobenius on \(A(-d)\) recorded after Definition \bpout{10.19}.


% object: sec-Local-terms-on-a-curve (blueprint/chap-The-trace-formula/sec-Local-terms-on-a-curve/_section.md), content hidden
\subsection*{\texorpdfstring{12.2\quad Local terms on a curve}{12.2 Local terms on a curve}}
\phantomsection\label{bp12}%
\addcontentsline{toc}{subsection}{\texorpdfstring{12.2\quad Local terms on a curve}{12.2 Local terms on a curve}}
% text: prose-12-06
\phantomsection\label{bp13}%
The local terms of the trace formula live on the fixed point scheme \(\Fix(\Phi_X^m)=\Fix(c_m)\) of the graph \(c_m\). More precisely, a local term is defined on each open and closed part of this scheme that is proper over \(k\). The next lemma shows that the fixed scheme is a finite reduced set of points, one for each \(\mathbf F_{q^m}\)-rational point of \(X_0\), so every such point has its own local term. For Frobenius on a smooth curve we then show that this local term is the trace on the stalk.

On \(\mathbf A^1\) the fixed scheme of \(\Phi^m\) is \(\operatorname{Spec}k[t]/(t^{q^m}-t)\). Its \(q^m\) points are reduced because \(t^{q^m}-t\) has derivative \(-1\). This is the algebraic form of the transversality seen at the start of Chapter \hyperref[bp1]{12}. Reducedness alone does not determine the local term, because the stalk does not determine \(\cHom(c_2^*K,c_1^!K)\) near the fixed point (Remark \hyperref[bp14]{12.4} below).

% object: remark-locality-warning (blueprint/chap-The-trace-formula/sec-Local-terms-on-a-curve/remark-locality-warning/_section.md)
\phantomsection\label{bp14}%
\begin{bpstatement}{Remark~12.4}{ (Local terms are not computed stalk by stalk)}
The vanishing of the local term at a rational point \(y\) with \(K_{\bar y}=0\), proved below, is a geometric fact about the restriction from cohomology with supports in the graph of \(\Phi_X^m\) to cohomology with supports in \(y\), and it does not follow from a locality property of internal Hom. For example, if \(j:\mathbf G_m\hookrightarrow\mathbf A^1\) and \(K=j_!A\), then the stalk of \(K\) at the origin is zero, but \(\cHom_A(j_!A,j_!A)\simeq Rj_*A\) has stalk \(A\) in degree zero there. Hence a local term cannot be computed by commuting a geometric stalk functor past this internal Hom.
\end{bpstatement}

% object: lemma-frobenius-fixed-scheme (blueprint/chap-The-trace-formula/sec-Local-terms-on-a-curve/lemma-frobenius-fixed-scheme/_section.md)
\phantomsection\label{bp15}%
\begin{bpstatementit}{Lemma~12.5}{ (The Frobenius fixed scheme)}
For \(m\geq1\), the fixed point scheme of \(\Phi_X^m\) satisfies

\[
\Fix(\Phi_X^m)
\simeq\coprod_{x\in X_0(\mathbf F_{q^m})}s.
\]

In particular, the fixed scheme is finite, reduced and proper over \(k\).
\end{bpstatementit}

% proof: proof/lemma-frobenius-fixed-scheme
\begin{proof}[{Proof}]\phantomsection\label{bp16}%
We reduce to a closed subscheme of affine space, where the equations \(t^{q^m}=t\) have simple roots.

For every open \(U_0\subset X_0\) the restriction to \(U=U_0\times_{k_0}k\) of the fixed scheme of \(\Phi_X^m\) is the fixed scheme of \(\Phi_U^m\), because \(\Phi_X\) is the base change of a morphism of \(X_0\) and so maps \(U\) into itself. As \(X_0\) is quasi-compact, we may therefore assume that \(X_0\) is a closed subscheme of \(\mathbf A^r_{\mathbf F_q}\), with coordinates \(t_1,\ldots,t_r\).

In this case the fixed scheme is the pullback of the diagonal along the graph. Since the equations of \(X_0\) have coefficients in \(\mathbf F_q\), we have \((\Phi_X^m)^\#t_i=t_i^{q^m}\), and the fixed scheme is cut out by the defining equations of \(X_0\) together with \(t_i^{q^m}-t_i=0\). Each polynomial \(t^{q^m}-t\) is separable, so the algebra \(k[t_1,\ldots,t_r]/(t_i^{q^m}-t_i)_i\) is a product of copies of \(k\) indexed by \(\mathbf F_{q^m}^r\). Its quotient by the defining equations of \(X_0\) retains exactly the factors at the rational solutions of those equations, as asserted. In particular the fixed scheme is a closed subscheme of a finite étale \(k\)-scheme, hence étale. The argument uses only the equations of \(X_0\), so \(X_0\) need not be reduced.
\end{proof}

% text: prose-12-07
\phantomsection\label{bp17}%
Thus every fixed point \(y\) is an open and closed point of the fixed scheme, and its local term \(\operatorname{LT}_y\) is defined.

% object: remark-outline-local-computation (blueprint/chap-The-trace-formula/sec-Local-terms-on-a-curve/remark-outline-local-computation/_section.md)
\phantomsection\label{bp18}%
\begin{bpstatement}{Remark~12.6}{ (Outline of the local computation)}
Let \(X_0\) be a smooth curve and \(y\in X_0(\mathbf F_{q^m})\). The trace formula on \(\mathbf P^1\) needs Proposition \hyperref[bp49]{12.16} below, which says that the local term at \(y\) of any morphism \(\psi:(\Phi_X^m)^*K\to K\) is the naive trace \(\operatorname{tr}_A(\psi_{\bar y}\mid K_{\bar y})\), also when \(K\) is not lisse at \(y\). To prove it, we filter \(K\) by the open complement of \(y\) and the point \(y\). By additivity of traces this reduces the proposition to two cases: the point itself, where Lemma \bpout{11.10} applies, and the case \(K_{\bar y}=0\), where we must show that the local term vanishes.

The second case needs a geometric argument. Indeed, the local term is defined on a neighbourhood of \((y,y)\) in \(X\times X\), where the Frobenius graph meets the diagonal, and the vanishing of the stalk at \(y\) does not by itself control it. For this we follow [\hyperref[bp105]{SGA5}, III B, §§2–4]. In this approach the local term is a restriction map from the graph to the diagonal, and blowing up \((y,y)\) separates the two curves. After the blowup the local term becomes a restriction to punctured curves near the point where two transverse divisors meet. There it vanishes when the monodromy is an \(\ell\)-group. In general we reduce to that case by passing to finite covers of degree prime to \(\ell\) (Appendix \bpout{D}); for such a cover the trace map splits the pullback.

What vanishes is a morphism in \(D(A)\), before any trace is taken. For this reason the vanishing holds for every \(\psi\), for every finite \(\ell\)-primary ring \(A\), and for every \(K\) with \(K_{\bar y}=0\). In particular it holds for \(j_!K\) on \(\mathbf P^1\), which is not lisse at infinity. We forget the \(A\)-action only in Lemma \hyperref[bp28]{12.9}, to detect a zero object.
\end{bpstatement}

% object: subsec-The-local-term-as-a-restriction-map (blueprint/chap-The-trace-formula/sec-Local-terms-on-a-curve/subsec-The-local-term-as-a-restriction-map/_section.md), content hidden
\subsubsection*{\texorpdfstring{12.2.1\quad The local term as a restriction map}{12.2.1 The local term as a restriction map}}
\phantomsection\label{bp19}%
\addcontentsline{toc}{subsubsection}{\texorpdfstring{12.2.1\quad The local term as a restriction map}{12.2.1 The local term as a restriction map}}
% text: prose-12-08
\phantomsection\label{bp20}%
The vanishing of a local term when \(K_{\bar y}=0\) is hard to see from the definition through the trace morphism (Definition \bpout{11.8}), so we first replace that definition by a description familiar from topology. There the intersection number of two submanifolds \(\Gamma\) and \(\Delta\) is computed by restricting the Thom class of \(\Gamma\) to \(\Delta\). Since the Thom class is supported on \(\Gamma\), the result is a class supported on \(\Gamma\cap\Delta\). Similarly, in our setting \(\psi\) defines a class supported on the graph. A restriction map \(\rho\) takes it to cohomology of the diagonal supported at \(y\), and the local term is the evaluation of the result. Note that the map \(\rho\) does not involve \(\psi\). It therefore suffices to show that \(\rho\) itself is zero when \(K_{\bar y}=0\) (Lemma \hyperref[bp45]{12.15}), and then the vanishing holds for every \(\psi\), invertible or not.

% object: lemma-local-term-restriction-map (blueprint/chap-The-trace-formula/sec-Local-terms-on-a-curve/subsec-The-local-term-as-a-restriction-map/lemma-local-term-restriction-map/_section.md)
\phantomsection\label{bp21}%
\begin{bpstatementit}{Lemma~12.7}{ (The local term as a restriction map)}
Let \(X_0\) be a smooth curve over \(k_0\), let \(X=X_0\times_{k_0}k\), let \(K\in D^b_{\mathrm{ctf}}(X,A)\), and let \(y\) be a fixed point of \(\Phi_X^m\). Let \(\Gamma\subset X\times X\) be the transposed graph of \(\Phi_X^m\), the image of the closed immersion \(\gamma=(\Phi_X^m,\operatorname{id}_X)\), and let \(\Delta\subset X\times X\) be the diagonal, the image of \(\delta=(\operatorname{id}_X,\operatorname{id}_X)\). On the strict localization \(T\) of \(X\times X\) at \((y,y)\) put

\[
P=(\mathbb D_XK\boxtimes K)|_T.
\]

A morphism \(\psi:(\Phi_X^m)^*K\to K\) determines a class \([\psi]\in H^0_\Gamma(T,P)\), and the local term \(\operatorname{LT}_y(\psi)\) of Definition \bpout{11.9} is the image of this class under the composite

\[
R\Gamma_\Gamma(T,P)
\xrightarrow{\rho}
R\Gamma_y(\Delta,\delta^*P)
\xrightarrow{\mathrm{ev}_K}
R\Gamma_y(X,\omega_X)
\simeq R\Gamma(y,\omega_y)
\xrightarrow{\int_y} A,
\]

where \(\rho\) is restriction from \(\Gamma\)-supported to \(y\)-supported cohomology.
\end{bpstatementit}

% text: prose-12-09
\phantomsection\label{bp22}%
In this lemma the two factors of the product are interchanged compared with the convention \(c_m=(\operatorname{id}_X,\Phi_X^m)\) of Definition \hyperref[bp7]{12.2}.

% proof: proof/lemma-local-term-restriction-map
\begin{proof}[{Proof of Lemma 12.7}]\phantomsection\label{bp23}%
We show that the local term and the displayed composite are both adjoint to one map \(R\gamma_*\gamma^!P\to R\delta_*\omega_X\). Once the class \([\psi]\) and the map \(\rho\) are built, the comparison is formal, and computing on étale models changes nothing.

By external duality, \(\gamma^!P\) is identified with \(\cHom((\Phi_X^m)^*K,K)\), so exceptional adjunction identifies \(\psi\) with \([\psi]\). The map \(\rho\) is constructed from the two canonical arrows

\[
R\gamma_*\gamma^!P\longrightarrow P\longrightarrow R\delta_*\delta^*P.
\]

After restriction to \(\Delta\), the source is supported on \(\Gamma\cap\Delta\). By Lemma \hyperref[bp15]{12.5} the scheme \(\Gamma\cap\Delta\) is reduced, and on \(T\) it is the single point \((y,y)\). Hence the resulting map factors through \(R\Gamma_y(\Delta,\delta^*P)\). Equivalently, it is the ordinary pullback of the object with supports, followed by the base change map to local cohomology with support in the intersection.

The adjoint of \(\mathrm{ev}_K:\delta^*P\to\omega_X\) is the map \(P\to R\delta_*\omega_X\) used to define the trace morphism in Definition \bpout{11.8}. Hence the displayed composite, before \(\int_y\), is adjoint to the map

\[
R\gamma_*\gamma^!P\longrightarrow P\longrightarrow R\delta_*\omega_X.
\]

The local term is built from the same map. Indeed, by exceptional proper base change for \(\delta\), the adjoint of this map, which is supported on \(y\), is the trace map \(\gamma^!P\to Ri_*\omega_y\) of Definition \bpout{11.9}. The triangle identities of the ordinary and exceptional adjunctions then show that the two constructions give the same element of \(A\).

The constructions above also make sense on the constructible étale models, where restriction to the strict localization is the filtered limit over étale neighbourhoods. The comparison theorems identify inverse images, exceptional adjunctions, duality and the support maps on the two sites, and hence they identify the class obtained on the étale models with the pro-étale local term of Definition \bpout{11.9}.
\end{proof}


% object: subsec-Vanishing-on-a-strictly-local-surface (blueprint/chap-The-trace-formula/sec-Local-terms-on-a-curve/subsec-Vanishing-on-a-strictly-local-surface/_section.md), content hidden
\subsubsection*{\texorpdfstring{12.2.2\quad Vanishing on a strictly local surface}{12.2.2 Vanishing on a strictly local surface}}
\phantomsection\label{bp24}%
\addcontentsline{toc}{subsubsection}{\texorpdfstring{12.2.2\quad Vanishing on a strictly local surface}{12.2.2 Vanishing on a strictly local surface}}
% text: prose-12-10
\phantomsection\label{bp25}%
The vanishing of the local term is explained by the position of the curves at \((y,y)\). Since \(d\Phi_X=0\), the transposed graph \(\Gamma\) is tangent to the vertical axis \(V=\{y\}\times X\). The diagonal \(\Delta\) and the horizontal axis \(H\) have other tangent directions. This persists on suitable covers of degree prime to \(\ell\) (Lemma \hyperref[bp40]{12.13} below). Blowing up the closed point separates the curves: the strict transforms of \(\Gamma\cup V\), of \(\Delta\) and of \(H\) meet the exceptional divisor \(E\simeq\mathbf P^1\) in different points.

Let \(M\) be the direct image of the coefficients from the complement of \(E\). We consider \(M\) on \(E\) minus the points of \(\Gamma\cup V\), and let \(z\) be the point where \(H\) meets \(E\). There \(M\) is locally constant away from \(z\), and its stalk at \(z\) is zero. Both facts use that the monodromy is an \(\ell\)-group (Lemma \hyperref[bp28]{12.9}). Now a branch of \(\Delta\) meets \(E\) in one point \(w\), and the restriction to this branch factors through the stalk \(M_w\). Composing the restriction from \(E\) to \(M_w\) with the specialization isomorphism from \(M_w\) to the generic stalk of \(E\), we get the restriction to the zero stalk \(M_z\) followed by specialization. Hence the restriction to \(M_w\) is zero.

The transverse case is Lemma \hyperref[bp28]{12.9}, and Lemma \hyperref[bp34]{12.11} reduces to it by blowing up the closed point. We compute on classical étale models of strictly local schemes, which are filtered limits of schemes of finite type, and in the proof of Lemma \hyperref[bp45]{12.15} we transport the result to the pro-étale site.

% object: lemma-punctured-surface-punctured-curve (blueprint/chap-The-trace-formula/sec-Local-terms-on-a-curve/subsec-Vanishing-on-a-strictly-local-surface/lemma-punctured-surface-punctured-curve/_section.md)
\phantomsection\label{bp26}%
\begin{bpstatementit}{Lemma~12.8}{ (A punctured surface and a punctured curve)}
Let \(S\) be the strict localization of a smooth surface over an algebraically closed field \(k\) of characteristic \(p\) at a closed point \(z\), let \(\ell\neq p\), and let \(D_1,D_2\subset S\) be regular divisors meeting transversely at \(z\). Then the restriction map

\[
R\Gamma(S-D_2,\mathbf F_\ell)\longrightarrow R\Gamma(D_1-z,\mathbf F_\ell)
\]

is an isomorphism.
\end{bpstatementit}

% proof: proof/lemma-punctured-surface-punctured-curve
\begin{proof}[{Proof}]\phantomsection\label{bp27}%
Both complexes have cohomology \(\mathbf F_\ell\) in degree zero, \(\mathbf F_\ell(-1)\) in degree one, and zero in all other degrees. For the surface this follows from localization, the acyclicity of the strictly local schemes \(S,D_2\), and the purity identity \(i_{D_2}^!\mathbf F_\ell=\mathbf F_\ell(-1)[-2]\). This identity, in turn, comes by transitivity from the dualizing complexes of the smooth schemes \(S\) and \(D_2\).

For the punctured trait \(D_1-z\), vanishing in degrees \(\geq2\) follows from Proposition \bpout{D.4}, by passing to the filtered limit of punctured neighbourhoods that are curves. In degree one the Kummer sequence gives \(H^1(-,\mu_\ell)=k'^*/k'^{*\ell}\), where \(k'\) is the fraction field. Moreover, the units of the strictly henselian trait are \(\ell\)-divisible. Hence the valuation induces \(k'^*/k'^{*\ell}\simeq\mathbf Z/\ell\), and the class of a uniformizer is a generator.

To see that the restriction map itself, and not only the two sides, is an isomorphism, we find a generator of \(H^1\) on each side. The restriction is clearly the identity in degree zero. In degree one, let \(g\) be a local equation of \(D_2\). Since the regular local ring of \(S\) is factorial, \(\operatorname{Pic}(S-D_2)=0\) and the units on \(S-D_2\) are the units of \(S\) times the powers of \(g\). As the units of \(S\) are \(\ell\)-divisible, the Kummer sequence shows that \(H^1(S-D_2,\mu_\ell)\simeq\mathbf Z/\ell\), generated by the class of \(g\). Since \(D_1\) and \(D_2\) meet transversely, \(g\) restricts to a uniformizer of \(D_1\), whose class generates \(H^1(D_1-z,\mu_\ell)\). So the restriction is an isomorphism in degree one as well, and since both complexes vanish in the other degrees, it is an isomorphism.
\end{proof}

% object: lemma-cohomology-near-two-transverse-divisors (blueprint/chap-The-trace-formula/sec-Local-terms-on-a-curve/subsec-Vanishing-on-a-strictly-local-surface/lemma-cohomology-near-two-transverse-divisors/_section.md)
\phantomsection\label{bp28}%
\begin{bpstatementit}{Lemma~12.9}{ (Cohomology near two transverse divisors)}
Let \(X\) be a smooth surface over an algebraically closed field of characteristic \(p\), or an étale limit with affine transition maps of such surfaces. Let \(D_1,D_2\) be regular divisors meeting transversely at exactly one point \(z\), and let

\[
j:X-D_2\hookrightarrow X,\qquad
j':X-(D_1\cup D_2)\hookrightarrow X-D_2.
\]

Let \(\ell\ne p\), let \(A\) be a finite commutative \(\ell\)-primary ring, and let \(N\) be a bounded complex of \(A\)-modules on \(X-(D_1\cup D_2)\) with finite locally constant cohomology whose monodromy factors through a finite \(\ell\)-group. Put \(M=Rj_*j'_!N\). Then \(M_z=0\), and \(M|_{D_2-z}\) has finite locally constant cohomology. If \(D_2\) is connected, then for every closed geometric point \(w\in D_2-z\) the restriction map

\[
R\Gamma(D_2,M|_{D_2})\longrightarrow M_w
\]

is zero in \(D(A)\).
\end{bpstatementit}

% proof: proof/lemma-cohomology-near-two-transverse-divisors
\begin{proof}[{Proof}]\phantomsection\label{bp29}%
The idea is that, since the monodromy is an \(\ell\)-group, we may filter the coefficients until they become the constant sheaf \(\mathbf F_\ell\). For constant coefficients, \(M_z=0\) is a statement on the strict localization of \(X\) at \(z\): restriction from the complement of \(D_2\) to the punctured curve \(D_1-z\) is an isomorphism on cohomology. This holds because a local equation of \(D_2\) gives the generator of \(H^1\) on both sides. Once we know that \(M\) is locally constant along \(D_2-z\), the vanishing of the restriction map follows from \(M_z=0\) by specialization along \(D_2\).

For the assertion at \(z\), pass to the strict localization \(S\) of \(X\) at \(z\) and write \(D_{1,S},D_{2,S}\) for the two transverse divisors. By the formula for stalks of a direct image along an open immersion [Stacks, \href{https://stacks.math.columbia.edu/tag/03Q7}{03Q7}] and localization we have

\[
M_z=R\Gamma(S-D_{2,S},j'_!N).
\]

To prove that this object is zero we may forget the \(A\)-action, since the exact forgetful functor is conservative and the cohomology of sheaves of modules agrees with the cohomology of the underlying abelian sheaves [Stacks, \href{https://stacks.math.columbia.edu/tag/03FD}{03FD}]. A finite abelian \(\ell\)-group with an action of a finite \(\ell\)-group has a finite filtration whose quotients are \(\mathbf F_\ell\) with trivial action. To see this, filter by powers of \(\ell\) and use the nilpotence of the augmentation ideal of the group algebra over \(\mathbf F_\ell\). Since \(j'_!\) and \(R\Gamma(S-D_{2,S},-)\) are exact, the class of \(N\) for which \(R\Gamma(S-D_{2,S},j'_!N)=0\) is closed under extensions. Using truncations and these composition series we reduce the assertion to \(N=\mathbf F_\ell\) on \(S-(D_{1,S}\cup D_{2,S})\).

By Lemma \hyperref[bp26]{12.8}, applied to \(S\) with the divisors \(D_{1,S},D_{2,S}\), the restriction map \(R\Gamma(S-D_{2,S},\mathbf F_\ell)\to R\Gamma(D_{1,S}-z,\mathbf F_\ell)\) is an isomorphism. The localization sequence for \(D_{1,S}-z\) and its complement in \(S-D_{2,S}\) identifies \(R\Gamma(S-D_{2,S},j'_!\mathbf F_\ell)\) with the cone of this map, shifted by \(-1\). Hence this complex is zero, and \(M_z=0\) in \(D(A)\) by conservativity.

We turn to the local constancy of \(M\) along \(D_2-z\), where only the divisor \(D_2\) remains. Since \(D_1\) meets \(D_2\) only at \(z\), we have \(j'_!N=N\) near \(D_2-z\), so \(M=Rj_*N\) there. To compute \(Rj_*N\) we use Abhyankar's lemma along the regular divisor \(D_2\) [Stacks, \href{https://stacks.math.columbia.edu/tag/0EYG}{0EYG}]: étale locally on \(D_2\), a finite tame local system is trivialized by a Kummer cover of a local parameter of \(D_2\) and an unramified cover. The lemma applies here because our finite \(\ell\)-group monodromy is tame, wild inertia being pro-\(p\). Consequently, the stalks of \(Rj_*N\) are the continuous cohomology of the same representation of tame inertia, and their specialization maps along \(D_2-z\) are isomorphisms. Explicitly, the pro-\(\ell\) part of tame inertia is \(\mathbf Z_\ell(1)\). Its cohomology with coefficients in a finite \(\ell\)-primary module is computed by the two-term complex given by a topological generator minus one. On the other hand, the finite quotients of the remaining factors of tame inertia have order prime to \(\ell\), so taking invariants under them is exact on these coefficients. Finally, we may choose the parameter and the unramified cover on an étale neighbourhood. Hence this description, including the specialization maps, is locally constant. Applying the argument to the cohomology sheaves of \(N\) and using the bounded spectral sequence proves the assertion for \(N\).

The last assertion combines the two facts just proved. A global section of \(M|_{D_2}\) specializes to every stalk. Its image at \(z\) is zero, and along \(D_2-z\) specialization is an isomorphism on stalks. To make this precise, choose a geometric generic point \(\bar\eta\) of the connected regular curve \(D_2\), together with specialization maps from \(\bar\eta\) to \(z\) and to \(w\). By naturality of restriction we get a commutative square

\[
\begin{CD}
R\Gamma(D_2,M|_{D_2}) @>>> M_w\\
@VVV @VV{\sim}V\\
M_z=0 @>>> M_{\bar\eta}.
\end{CD}
\]

The right arrow is an isomorphism by the local constancy just proved, so the upper arrow is zero.
\end{proof}

% text: prose-12-11
\phantomsection\label{bp30}%
Since the restriction to a curve meeting \(D_2\) in one point factors through the stalk of \(M\) at that point, the lemma implies that restrictions to punctured curves vanish.

% object: lemma-restriction-punctured-curve (blueprint/chap-The-trace-formula/sec-Local-terms-on-a-curve/subsec-Vanishing-on-a-strictly-local-surface/lemma-restriction-punctured-curve/_section.md)
\phantomsection\label{bp31}%
\begin{bpstatementit}{Lemma~12.10}{ (Restriction to a punctured curve)}
In the situation of Lemma \hyperref[bp28]{12.9} with \(D_2\) connected, let \(C\subset X\) be a closed curve with \(C\cap D_2=\{w\}\), and let \(C_{(w)}\) be its strict localization at \(w\). Then the restriction map

\[
R\Gamma(X-D_2,j'_!N)\longrightarrow
R\Gamma(C_{(w)}-w,(j'_!N)|_{C_{(w)}-w})
\]

is zero in \(D(A)\).
\end{bpstatementit}

% proof: proof/lemma-restriction-punctured-curve
\begin{proof}[{Proof}]\phantomsection\label{bp32}%
The source is \(R\Gamma(X-D_2,j'_!N)=R\Gamma(X,M)\), by the definition \(M=Rj_*j'_!N\). Since \(C\) meets \(D_2\) only at \(w\), the punctured curve \(C_{(w)}-w\) lies over the strict localization of \(X\) at \(w\) and avoids \(D_2\). The map therefore factors as the restriction to \(R\Gamma(D_2,M|_{D_2})\), followed by the restriction to the stalk \(M_w=R\Gamma(X_{(w)}-D_2,j'_!N)\) (by the stalk formula), followed by the restriction to the displayed target. Since all these maps are restrictions of the same complex, the composite is the given map. If \(w\ne z\), the middle map is zero by Lemma \hyperref[bp28]{12.9}. If \(w=z\), the intermediate object \(M_z\) is zero.
\end{proof}

% text: prose-12-12
\phantomsection\label{bp33}%
Near \((y,y)\) the curves we care about do not meet transversely, since the graph, the diagonal and the axis \(H\) all pass through the same point. The next lemma reduces this configuration to the transverse one by a single blowup.

% object: lemma-vanishing-restrictions-after-blowing-up (blueprint/chap-The-trace-formula/sec-Local-terms-on-a-curve/subsec-Vanishing-on-a-strictly-local-surface/lemma-vanishing-restrictions-after-blowing-up/_section.md)
\phantomsection\label{bp34}%
\begin{bpstatementit}{Lemma~12.11}{ (Vanishing of restrictions after blowing up)}
Let \(T\) be the strict localization of a smooth surface at a closed point \(t\). Let \(H\) be a regular curve through \(t\), and let \(D\) and \(C\) be reduced closed curves through \(t\) such that the tangent cone of \(D\) at \(t\) meets the tangent cones of \(C\) and of \(H\) only in \(0\). Let \(j':T-(D\cup H)\hookrightarrow T-D\) be the open immersion and let \(N\) be a bounded complex on \(T-(D\cup H)\) with finite locally constant cohomology and finite \(\ell\)-group monodromy. Then the restriction map

\[
R\Gamma(T-D,j'_!N)\longrightarrow
R\Gamma(C-t,(j'_!N)|_{C-t})
\]

is zero in \(D(A)\).
\end{bpstatementit}

% proof: proof/lemma-vanishing-restrictions-after-blowing-up
\begin{proof}[{Proof}]\phantomsection\label{bp35}%
After blowing up \(t\), the three curves are separated along the exceptional divisor, which then plays the role of \(D_2\) in the previous lemma.

Let \(\widetilde T\to T\) be the blowup in \(t\), with exceptional divisor \(E\simeq\mathbf P^1_k\), and let \(D',H',C'\) be the strict transforms. In these terms, the hypotheses on tangent cones say that \(D'\cap E\) is disjoint from both \(C'\cap E\) and \(H'\cap E\). We therefore apply Lemmas \hyperref[bp28]{12.9} and \hyperref[bp31]{Lemma 12.10} to the surface \(X=\widetilde T-D'\) with the divisors \(D_1=H'\cap X\) and \(D_2=E\cap X=E-D'\). Here the divisor \(E-D'\) is connected, since it is a projective line with finitely many points removed, and \(H'\cap X\) meets it transversely in the single point \(z\) corresponding to the tangent direction of \(H\). Since the blowup is an isomorphism off \(E\), the source of the map is \(R\Gamma(X-E,j'_!N')\) for the pulled-back complex \(N'\).

The strict transform \(C'\) is proper and quasi-finite over the strictly local curve \(C\), hence finite. Its closed fibre is the finite set of tangent directions of its branches, so \(C'\) is a finite disjoint union of strictly local curves \(C'_w\), each meeting \(E\) at a single point \(w\) of \(E-D'\). Each restriction to \(C'_w-w\) is zero by Lemma \hyperref[bp31]{12.10}. Since \(C'-E\simeq C-t\), the given restriction map is the direct sum of these zero maps.
\end{proof}


% object: subsec-The-local-term-at-a-zero-stalk (blueprint/chap-The-trace-formula/sec-Local-terms-on-a-curve/subsec-The-local-term-at-a-zero-stalk/_section.md), content hidden
\subsubsection*{\texorpdfstring{12.2.3\quad The local term at a zero stalk}{12.2.3 The local term at a zero stalk}}
\phantomsection\label{bp36}%
\addcontentsline{toc}{subsubsection}{\texorpdfstring{12.2.3\quad The local term at a zero stalk}{12.2.3 The local term at a zero stalk}}
% text: prose-12-13
\phantomsection\label{bp37}%
In the setting of Lemma \hyperref[bp21]{12.7} we compute on the classical étale models. This is legitimate because the classical comparison identifies tensor products, duality, exceptional adjunctions and the pairing between graph and diagonal with their versions there. Let \(X_1\) and \(X_2\) be two copies of the strict localization of \(X\) at \(y\). The strict localization of \(X_1\times X_2\) at \((y,y)\) is the strictly local surface \(T\) of Lemma \hyperref[bp21]{12.7}. Write

\[
H=X_1\times\{y\},\qquad V=\{y\}\times X_2,
\]

and keep \(\Gamma\) and \(\Delta\) for the transposed Frobenius graph and the diagonal in \(T\). Let \(u\) and \(v\) be the pullbacks to the two factors \(X_1\) and \(X_2\) of the same uniformizer of \(X\) at \(y\), so that \(\Delta\), \(V\) and \(H\) have equations \(u-v=0\), \(u=0\) and \(v=0\). Write \(\varphi(v)=(\Phi_X^m)^\#u\), a power series in \(v\); the graph \(\Gamma\) has equation \(u-\varphi(v)=0\). Since \(d\Phi_X=0\), the graph \(\Gamma\) is tangent to \(V\), while \(\Delta\) and \(H\) have other tangent directions.

The next three lemmas prepare the proof of Lemma \hyperref[bp45]{12.15}.

% object: lemma-covers-degree-prime-to-ell (blueprint/chap-The-trace-formula/sec-Local-terms-on-a-curve/subsec-The-local-term-at-a-zero-stalk/lemma-covers-degree-prime-to-ell/_section.md)
\phantomsection\label{bp38}%
\begin{bpstatementit}{Lemma~12.12}{ (Covers of degree prime to $\ell$)}
Let \(S\) be the strict localization of a smooth curve over the algebraically closed field \(k\) at a closed point \(s\), and let \(\mathcal F_1,\dots,\mathcal F_r\) be finite locally constant sheaves on \((S-s)_{\acute{e}t}\). There is an integer \(d_1\) prime to \(\ell\) with the following property. For every multiple \(d\) of \(d_1\) that is prime to \(\ell\) there is a finite flat morphism \(S'\to S\) of degree \(d\) such that

\begin{enumerate}
\item \(S'\) is the strict localization of a smooth curve over \(k\) at a closed point \(s'\), and \(s'\) is the only point of \(S'\) over \(s\); more precisely, \(S'\to S\) is the base change of a finite flat morphism \(C'\to C\) of smooth curves over \(k\), where \(C\) is an étale neighbourhood of \(s\), which is étale away from \(s\) and has a single point over \(s\),
\item \(S'-s'\to S-s\) is finite étale, and
\item the monodromy of the pullback of each \(\mathcal F_i\) to \(S'-s'\) factors through a finite \(\ell\)-group.
\end{enumerate}
\end{bpstatementit}

% proof: proof/lemma-covers-degree-prime-to-ell
\begin{proof}[{Proof}]\phantomsection\label{bp39}%
We take the fixed field of an \(\ell\)-Sylow subgroup of the monodromy, adjust the degree by an Eisenstein extension, and descend to curves.

Let \(F\) be the fraction field of the discrete valuation ring of \(S\). Then the punctured trait \(S-s\) is \(\operatorname{Spec}(F)\), so the \(\mathcal F_i\) are finite sets with a continuous action of \(\operatorname{Gal}(F^{\mathrm{sep}}/F)\). We choose a finite Galois extension \(L/F\) over which all of them become constant, let \(P\) be an \(\ell\)-Sylow subgroup of \(G=\operatorname{Gal}(L/F)\), and put \(F_1=L^P\) and \(d_1=[F_1:F]=[G:P]\), which is prime to \(\ell\).

Let \(d=d_1e\) be a multiple of \(d_1\) prime to \(\ell\). The normalization \(S_1\) of \(S\) in the separable extension \(F_1\) is finite over \(S\), and it is local because it is a domain finite over a henselian local ring. Hence \(S_1\) is a strictly henselian discrete valuation ring with residue field \(k\), and it has degree \(d_1\) over \(S\). If \(e>1\), let \(t\) be a uniformizer of \(S_1\). The Eisenstein polynomial \(z^e+tz-t\) is irreducible, and its derivative \(ez^{e-1}+t\) is nonzero, also when \(p\mid e\), so it is separable. Hence \(S'=S_1[z]/(z^e+tz-t)\) is a discrete valuation ring, totally ramified of degree \(e\) over \(S_1\), with separable fraction field \(F'\) over \(F\). If \(e=1\) we put \(S'=S_1\). Then \(S'\to S\) is finite flat of degree \(d\), with a single point \(s'\) over \(s\), and \(S'-s'=\operatorname{Spec}(F')\to
\operatorname{Spec}(F)\) is finite étale. The group \(\operatorname{Gal}(F^{\mathrm{sep}}/F')\) acts on the stalks of the \(\mathcal F_i\) through its image in \(G\), which lies in \(\operatorname{Gal}(L/F_1)=P\) since \(F_1\subset F'\).

Finally, by a limit argument \(S'\to S\) descends to a finite flat morphism \(C'\to C\) onto an étale neighbourhood \(C\) of \(s\), which we may take with normal source. Since \(k\) is perfect, the normal curve \(C'\) is smooth. Moreover, the morphism \(C'\to C\) is étale over a dense open of \(C\), and after shrinking \(C\) around \(s\) it is étale away from \(s\). Since \(k\) is algebraically closed, \(C'\times_CS\) is the product of the strict localizations of \(C'\) at the points over \(s\). As it equals the local scheme \(S'\), there is a single point over \(s\), and \(S'\) is the strict localization of \(C'\) there.
\end{proof}

% object: lemma-tangent-cones-on-cover (blueprint/chap-The-trace-formula/sec-Local-terms-on-a-curve/subsec-The-local-term-at-a-zero-stalk/lemma-tangent-cones-on-cover/_section.md)
\phantomsection\label{bp40}%
\begin{bpstatementit}{Lemma~12.13}{ (Tangent cones on the cover)}
For \(i=1,2\) let \(X_i'\to X_i\) be a finite flat morphism of degree \(d\) as in Lemma \hyperref[bp38]{12.12}, with \(y_i'\) the point over \(y\). Let \(T'\) be the strict localization of \(X_1'\times X_2'\) at \(y'=(y_1',y_2')\), and write \(\Gamma',V',\Delta',H'\) for the inverse images of \(\Gamma,V,\Delta,H\) in \(T'\). Then \(T'\) is the strict localization of a smooth surface over \(k\) at a closed point, the morphism \(T'\to T\) is finite flat of degree \(d^2\) and étale over \(T-(H\cup V)\), the reduction \(H'_{\mathrm{red}}\) is regular, and the tangent cone of \((\Gamma'\cup V')_{\mathrm{red}}\) at \(y'\) meets the tangent cones of \(\Delta'_{\mathrm{red}}\) and of \(H'_{\mathrm{red}}\) only in \(0\).
\end{bpstatementit}

% proof: proof/lemma-tangent-cones-on-cover
\begin{proof}[{Proof}]\phantomsection\label{bp41}%
The two covers descend to finite flat morphisms \(C_i'\to C_i\) of degree \(d\) between smooth curves, where \(C_i\) is an étale neighbourhood of \(y\) in \(X\), étale away from \(y\) and with a single point \(y_i'\) over \(y\) (Lemma lemma-covers-degree-prime-to-ell). Their product \(C_1'\times C_2'\to C_1\times C_2\) is a finite flat morphism of degree \(d^2\) of smooth surfaces, étale away from the two fibres through \(y\), with \(y'\) the only point over \((y,y)\). Base change to the strict localization \(T\) of \(C_1\times C_2\) at \((y,y)\) gives a scheme finite over \(T\) with one point over the closed point, hence local and strictly henselian, and this is \(T'\). This proves the assertions about \(T'\) and \(T'\to T\).

Choose uniformizers \(\xi\) of \(X_1'\) and \(\eta\) of \(X_2'\), so that the completed local ring of \(T'\) at \(y'\) is \(k[[\xi,\eta]]\). Since the residue field \(k\) is algebraically closed, both covers are totally ramified, so \(u\) is a unit of \(k[[\xi]]\) times \(\xi^d\), and similarly for \(v\). Multiplying \(\xi\) and \(\eta\) by suitable \(d\)th roots of nonzero constants, we may assume that

\[
u\equiv\xi^d,\qquad v\equiv\eta^d\pmod{(\xi,\eta)^{d+1}}.
\]

Since \(\Phi_X^m\) is purely inseparable of degree \(q^m\) and fixes \(y\), we have \(\operatorname{ord}_v\varphi=q^m>1\), and hence \(\varphi(v)\in(\eta)^{dq^m}\subseteq(\xi,\eta)^{d+1}\). It follows that the pulled-back equations of \(\Gamma'\), \(V'\), \(\Delta'\) and \(H'\) have initial forms

\[
\xi^d,\qquad\xi^d,\qquad\xi^d-\eta^d,\qquad\eta^d,
\]

respectively. The support of the tangent cone of a curve is the zero set of the initial form of its equation, and it does not change when we pass to the reduction, since if \(J\) is the radical of \(I\) and \(J^n\subset I\) then the ideals of initial forms satisfy \(\mathrm{in}(J)^n\subset\mathrm{in}(I)
\subset\mathrm{in}(J)\). Thus the tangent cone of \((\Gamma'\cup V')_{\mathrm{red}}\) is supported on the line \(\xi=0\), that of \(H'_{\mathrm{red}}\) on \(\eta=0\), and that of \(\Delta'_{\mathrm{red}}\) on the lines \(\xi=\zeta\eta\) with \(\zeta^d=1\). This holds also when \(p\mid d\) and \(\xi^d-\eta^d\) is a \(p\)th power. Away from the origin the first support meets neither of the others. Finally \(H'_{\mathrm{red}}\) is the curve \(\eta=0\), which is regular.
\end{proof}

% object: lemma-descent-vanishing-finite-etale (blueprint/chap-The-trace-formula/sec-Local-terms-on-a-curve/subsec-The-local-term-at-a-zero-stalk/lemma-descent-vanishing-finite-etale/_section.md)
\phantomsection\label{bp42}%
\begin{bpstatementit}{Lemma~12.14}{ (Descent of vanishing along a finite étale cover)}
Let \(f:W'\to W\) be a finite étale morphism of schemes of constant degree \(n\), with \(n\) invertible in \(A\), and let \(L\in D(W_{\acute{e}t},A)\). Then the pullback map \(R\Gamma(W,L)\to R\Gamma(W',f^{-1}L)\) has a left inverse in \(D(A)\). In particular, a morphism \(\alpha:M\to R\Gamma(W,L)\) in \(D(A)\) is zero if its composite with this pullback map is zero.
\end{bpstatementit}

% proof: proof/lemma-descent-vanishing-finite-etale
\begin{proof}[{Proof}]\phantomsection\label{bp43}%
Since \(f\) is finite étale, \(f_*=f_!\) is exact, and the trace map \(f_*f^{-1}\mathcal F\to\mathcal F\) composed with the adjunction unit \(\mathcal F\to f_*f^{-1}\mathcal F\) is multiplication by \(n\) for every sheaf of \(A\)-modules \(\mathcal F\) [Stacks, \href{https://stacks.math.columbia.edu/tag/03SH}{03SH}]. Since both maps are natural and the functors are exact, applying them termwise to a complex representing \(L\) gives maps \(L\to f_*f^{-1}L\to L\) in \(D(W_{\acute{e}t},A)\) whose composite is \(n\). Applying \(R\Gamma(W,-)\) and using \(R\Gamma(W,f_*f^{-1}L)=R\Gamma(W',f^{-1}L)\), we see that \(n^{-1}\) times the map induced by the trace is a left inverse of the pullback.
\end{proof}

% text: prose-12-14
\phantomsection\label{bp44}%
The map \(\rho\) of Lemma \hyperref[bp21]{12.7} depends on \(K\) and on the geometry of \(\Gamma\) and \(\Delta\) near \((y,y)\), but not on \(\psi\). We therefore state the vanishing of local terms at a zero stalk as a statement about \(\rho\) alone.

% object: lemma-vanishing-local-term-zero-stalk (blueprint/chap-The-trace-formula/sec-Local-terms-on-a-curve/subsec-The-local-term-at-a-zero-stalk/lemma-vanishing-local-term-zero-stalk/_section.md)
\phantomsection\label{bp45}%
\begin{bpstatementit}{Lemma~12.15}{ (Vanishing of the local term at a zero stalk)}
In the setting of Lemma \hyperref[bp21]{12.7}, let \(y\) be a fixed point of \(\Phi_X^m\), and suppose \(K_{\bar y}=0\). Then

\[
\rho:R\Gamma_\Gamma(T,P)\longrightarrow R\Gamma_y(\Delta,\delta^*P)
\]

is zero in \(D(A)\). Consequently \(\operatorname{LT}_y(\psi)=0\) for every morphism \(\psi:(\Phi_X^m)^*K\to K\).
\end{bpstatementit}

% proof: proof/lemma-vanishing-local-term-zero-stalk
\begin{proof}[{Proof}]\phantomsection\label{bp46}%
We use the notation introduced before Lemma \hyperref[bp38]{12.12} and compute on the étale model: we identify \(\rho\) with a restriction map \(r\) away from \(y\), kill \(r\) on a cover with \(\ell\)-group monodromy, and descend.

\emph{Step 1: reduction to \(r\).} Let \(N=(\mathbb D_XK)|_{X_1-y}\), the dual of \(K|_{X_1-y}\) with respect to the restriction of \(\omega_X\), and let \(Q=K|_{X_2-y}\). Write \(j_1:X_1-y\hookrightarrow X_1\) and \(j_2:X_2-y\hookrightarrow X_2\) for the inclusions. As \(K_{\bar y}=0\), localization gives \(K|_{X_2}=j_{2!}Q\), and \((\mathbb D_XK)|_{X_1}=Rj_{1*}N\) because the dual of an extension by zero is the direct image of the dual. The ordinary external Künneth comparison therefore identifies \(P=(\mathbb D_XK\boxtimes K)|_T\) with \(Rj_*j'_!(N\boxtimes Q)\), where

\[
j:T-V\hookrightarrow T,\qquad
j':T-(H\cup V)\hookrightarrow T-V.
\]

Both \(R\Gamma(T,P)\) and \(R\Gamma(\Delta,\delta^*P)\) are zero, because the schemes are strictly local and the stalk at the closed point is \((\mathbb D_XK)_{\bar y}\otimes^{\mathbf L} K_{\bar y}=0\). So the localization triangles identify \(\rho\), shifted by one, with the ordinary restriction \(R\Gamma(T-\Gamma,P)\to R\Gamma(\Delta-y,\delta^*P)\), which under the Künneth identification is

\[
r:R\Gamma(T-(V\cup\Gamma),j'_!(N\boxtimes Q))
\longrightarrow R\Gamma(\Delta-y,(N\boxtimes Q)|_{\Delta-y}).
\]

\emph{Step 2: the restriction vanishes on a cover.} Choose bounded finite projective lisse models of \(N\) and \(Q\) by Lemma \bpout{3.13}; their terms are finite locally constant, as a punctured strictly local trait is its generic point. Then Lemma \hyperref[bp38]{12.12}, applied to the terms of the model of \(N\) on \(X_1-y\) and to those of \(Q\) on \(X_2-y\), gives integers \(d_1,d_2\) prime to \(\ell\). With \(d=\operatorname{lcm}(d_1,d_2)\) we then obtain covers \(X_1'\to X_1\) and \(X_2'\to X_2\) of the same degree \(d\). We want equal degrees so that \(u\) and \(v\) become \(d\)th powers for the same \(d\) in Lemma \hyperref[bp40]{12.13}. Let \(T'\), \(\Gamma'\), \(V'\), \(\Delta'\), \(H'\) be as in that lemma, and let \(N'\) and \(Q'\) be the pullbacks of \(N\) and \(Q\) to \(X_1'-y_1'\) and \(X_2'-y_2'\). By construction, the complex \(N'\boxtimes Q'\) has finite \(\ell\)-group monodromy on \(T'-(H'\cup V')\), hence also on the smaller open set obtained by removing \(\Gamma'\). Then the restriction \(r'\) on \(T'\), defined like \(r\), is zero by Lemma \hyperref[bp34]{12.11}, applied to \(T'\) with \(D=(\Gamma'\cup V')_{\mathrm{red}}\), the regular curve \(H'_{\mathrm{red}}\) and \(C=\Delta'_{\mathrm{red}}\); its hypotheses hold by Lemma \hyperref[bp40]{12.13}, and passing to reductions changes neither the open complements nor the étale cohomology.

\emph{Step 3: descent.} The map \(r\) followed by pullback to \(\Delta'-y'\) equals pullback to \(T'\) followed by \(r'=0\), because pullback along \(T'\to T\) commutes with restriction and with extension by zero. As \(\Delta\) meets \(H\cup V\) only at \(y\), Lemma \hyperref[bp40]{12.13} shows that \(\Delta'-y'\to\Delta-y\) is finite étale of degree \(d^2\). This degree is invertible in \(A\), so Lemma \hyperref[bp42]{12.14} gives \(r=0\), and hence \(\rho=0\). By the comparison theorems, this identity on the classical étale models of these strictly local schemes also holds in the pro-étale formulation. Finally, by Lemma \hyperref[bp21]{12.7}, \(\operatorname{LT}_y(\psi)\) is the image of \([\psi]\) under a composite that begins with \(\rho\), so it is zero.
\end{proof}


% object: subsec-The-local-term-at-a-fixed-point (blueprint/chap-The-trace-formula/sec-Local-terms-on-a-curve/subsec-The-local-term-at-a-fixed-point/_section.md), content hidden
\subsubsection*{\texorpdfstring{12.2.4\quad The local term at a fixed point}{12.2.4 The local term at a fixed point}}
\phantomsection\label{bp47}%
\addcontentsline{toc}{subsubsection}{\texorpdfstring{12.2.4\quad The local term at a fixed point}{12.2.4 The local term at a fixed point}}
% text: prose-12-15
\phantomsection\label{bp48}%
The point \(y\) is stable under \(\Phi_X^m\), so additivity of trace classes applies to the localization triangle for \(X\setminus\{y\}\) and \(\{y\}\). Here the open piece has zero stalk at \(y\), so its local term vanishes by Lemma \hyperref[bp45]{12.15}. The closed piece, in turn, is pushed forward from the point, where the local term is a trace by Lemma \bpout{11.10}. Note that the argument never uses that \(K\) is lisse at \(y\). This is why, in the proof of Theorem \hyperref[bp62]{12.20}, we can apply the result to \(j_!K\) at the point at infinity of \(\mathbf P^1\).

% object: proposition-frobenius-local-term-curve (blueprint/chap-The-trace-formula/sec-Local-terms-on-a-curve/subsec-The-local-term-at-a-fixed-point/proposition-frobenius-local-term-curve/_section.md)
\phantomsection\label{bp49}%
\begin{bpstatementit}{Proposition~12.16}{ (The Frobenius local term on a curve)}
Let \(X_0\) be a smooth curve over \(k_0\), let \(X=X_0\times_{k_0}k\), let \(K\in D^b_{\mathrm{ctf}}(X,A)\) and let \(\psi:(\Phi_X^m)^*K\to K\) be any morphism. For \(y\in X_0(\mathbf F_{q^m})\) one has

\[
\operatorname{LT}_y(\psi)
=\operatorname{tr}_A(\psi_{\bar y}\mid K_{\bar y}).
\]
\end{bpstatementit}

% proof: proof/proposition-frobenius-local-term-curve
\begin{proof}[{Proof}]\phantomsection\label{bp50}%
Let \(i:\{y\}\hookrightarrow X\) and \(j:U=X\setminus\{y\}\hookrightarrow X\). We split \(\psi\) along the localization triangle and compute the two local terms separately.

\emph{Splitting \(\psi\).} Both pieces are stable under \(\Phi_X^m\), and in fact \(((\Phi_X^m)^{-1}\{y\})_{\mathrm{red}}=\{y\}\). In the enhanced derived category, the functorial localization triangle is a two-step filtered object,

\[
K_U=j_!j^*K\longrightarrow K\longrightarrow i_*i^*K
\xrightarrow{+1},
\]

and \(\psi\) induces a coherent map from the \((\Phi_X^m)^*\)-pullback of this filtration to the filtration itself, with induced correspondences \(\psi_U\) and \(i_*\psi_{\bar y}\) on the pieces. This rests on two base change identities. For the open part, \((\Phi_X^m)^*j_!\simeq j_!(\Phi_X^m|_U)^*\) by base change along inverse images in the cartesian square. The square for the point, by contrast, is cartesian only after passing to reductions, so there proper base change followed by nilpotent invariance gives \((\Phi_X^m)^*i_*\simeq i_*(\Phi_X^m|_y)^*\). Both comparisons are mates of the adjunction units for \(j\) and \(i\) used in the triangle, which is why the induced maps are compatible with the filtration.

\emph{Additivity.} The second half of Proposition \bpout{11.19} applies to the graph \(c_m\), because its hypothesis \(\Phi_X^m(\{y\})\subset\{y\}\) holds for the fixed point \(y\). It gives

\[
\operatorname{LT}_y(\psi)
=\operatorname{LT}_y(\psi_U)
+\operatorname{LT}_y(i_*\psi_{\bar y}).
\]

\emph{The two local terms.} The first term is zero by Lemma \hyperref[bp45]{12.15}, since \((K_U)_{\bar y}=0\). For the second, the description of the closed piece \(u_Z\) in Proposition \bpout{11.19} shows that \(i_*\psi_{\bar y}\) is the pushforward along \(i\) of the cohomological correspondence \(\psi_{\bar y}\) on \((y,K_{\bar y})\) for the identity correspondence of \(y\). Its local term is therefore \(\operatorname{tr}_A(\psi_{\bar y}\mid K_{\bar y})\), the trace on the perfect complex \(K_{\bar y}\), by the compatibility with proper pushforward in Theorem \bpout{11.13} and Lemma \bpout{11.10}.
\end{proof}

% text: prose-12-16
\phantomsection\label{bp51}%
For \(K=A_X\) on \(X=\mathbf A^1_k\) the proposition gives \(\operatorname{LT}_y(\psi_m)=1\) at each of the \(q^m\) fixed points of \(\Phi^m_X\). This matches the intersection count at the start of the chapter. For general \(K\) the proposition is not a formal consequence of the definitions, as Remark \hyperref[bp14]{12.4} explains.

% object: example-frobenius-projective-line (blueprint/chap-The-trace-formula/sec-Local-terms-on-a-curve/subsec-The-local-term-at-a-fixed-point/example-frobenius-projective-line/_section.md)
\phantomsection\label{bp52}%
\begin{bpstatement}{Example~12.17}{ (Frobenius on the projective line)}
Let \(X=\mathbf P^1_k\), \(K=A_X\), and let \(\psi\) be the constant Frobenius correspondence, so that \(\psi_m\) is the identity on every stalk. Both sides of the Lefschetz–Verdier formula for the graph \(c_m\) can be computed explicitly. The fixed points of \(\Phi_X^m\) are the \(q^m+1\) points of \(\mathbf P^1(\mathbf F_{q^m})\). By the proposition each has local term \(\operatorname{tr}_A(1\mid A)=1\), so the local side is \(q^m+1\). On the global side, \(R\Gamma(\mathbf P^1,A)=A\oplus A(-1)[-2]\). By Lemma \hyperref[bp9]{12.3} the operator \(T_{c_m}(\psi_m)\) is the \(m\)-th power of geometric Frobenius. It acts by \(1\) on \(H^0\) and by \(q^m\) on \(H^2=A(-1)\), and both degrees are even, so its trace is \(1+q^m\) with no signs. Compare the identity correspondence at the start of Section \bpout{11.4}. There the fixed locus is all of \(\mathbf P^1\) and the answer \(2\) is a self-intersection. Here the fixed locus is finite and each point counts once.
\end{bpstatement}



% object: sec-Frobenius-and-Rf (blueprint/chap-The-trace-formula/sec-Frobenius-and-Rf/_section.md), content hidden
\subsection*{\texorpdfstring{12.3\quad Frobenius and $Rf_!$}{12.3 Frobenius and Rf!}}
\phantomsection\label{bp53}%
\addcontentsline{toc}{subsection}{\texorpdfstring{12.3\quad Frobenius and $Rf_!$}{12.3 Frobenius and Rf!}}
% text: prose-12-17
\phantomsection\label{bp54}%
The trace formula with finite coefficients is proved by induction on the number of affine coordinates. We project a closed subscheme of \(\mathbf A^d\) to \(\mathbf A^{d-1}\), apply the induction hypothesis to \(Rf_!K\), and identify each stalk of \(Rf_!K\) at a rational point with the cohomology of the fibre, a closed subscheme of \(\mathbf A^1\) over a larger finite field. For this, \(Rf_!K\) must carry a Frobenius correspondence with two properties: its operator on \(R\Gamma_c\) is that of \(\psi\), and its stalk maps are the Frobenius operators on the cohomology of the fibres.

The square formed by \(f\) and the two Frobenii commutes, but it is usually not cartesian. Indeed, already for \(f:\mathbf A^1\to s\) the Frobenius of \(s\) is the identity, so the fibre product is \(\mathbf A^1\) itself, while \(\Phi_{\mathbf A^1}\) has degree \(q\). The difference is measured by the relative Frobenius. It is a universal homeomorphism, so pullback along it is an equivalence on étale sheaves, and the comparison map built from it is an isomorphism.

% object: lemma-comparison-via-relative-frobenius (blueprint/chap-The-trace-formula/sec-Frobenius-and-Rf/lemma-comparison-via-relative-frobenius/_section.md)
\phantomsection\label{bp55}%
\begin{bpstatementit}{Lemma~12.18}{ (Comparison via relative Frobenius)}
Let \(f_0:X_0\to Y_0\) be a morphism of separated schemes of finite type over \(k_0\), write \(f:X\to Y\) for its geometric base change, and let \(m\geq1\). Put

\[
X'_m=Y\times_{\Phi_Y^m,Y,f}X,
\quad p_m:X'_m\to X,\quad f'_m:X'_m\to Y,
\quad r_m=(f,\Phi_X^m):X\to X'_m,
\]

where \(p_m\) and \(f'_m\) are the two projections.

\begin{enumerate}
\item The relative Frobenius \(r_m\) is a finite universal homeomorphism, \(p_mr_m=\Phi_X^m\) and \(f'_mr_m=f\), and there are natural isomorphisms \[
   \begin{aligned}
   \theta_{f,m}:(\Phi_Y^m)^*Rf_!K
   &\simeq Rf'_{m!}p_m^*K
   \xrightarrow{Rf'_{m!}\eta_{r_m}}
   Rf'_{m!}Rr_{m*}r_m^*p_m^*K\\
   &\simeq Rf_!(\Phi_X^m)^*K,
   \end{aligned}
   \] compatible with iteration in \(m\).
\item For a Frobenius correspondence \(\psi:\Phi_X^*K\to K\) on \(K\in D^b_{\mathrm{ctf}}(X,A)\), the object \(L=Rf_!K\) of \(D^b_{\mathrm{ctf}}(Y,A)\) carries the Frobenius correspondence \(\psi'=Rf_!\psi\circ\theta_{f,1}:\Phi_Y^*L\to L\), and its iterates are \[
   \psi'_m=Rf_!\psi_m\circ\theta_{f,m}.
   \]
\item The composition isomorphism \(R\Gamma_c(Y,L)\simeq R\Gamma_c(X,K)\) intertwines the endomorphisms \(T_{c_m}(\psi'_m)\) on \(Y\) and \(T_{c_m}(\psi_m)\) on \(X\) of Definition \hyperref[bp7]{12.2}.
\item At \(y\in Y_0(\mathbf F_{q^m})\), base change for compact supports identifies \((L_{\bar y},\psi'_{m,\bar y})\) with \[
   \left(R\Gamma_c(X_{\bar y},K|_{X_{\bar y}}),
   T_{c_1}(\psi_m|_{X_{\bar y}})\right),
   \] where the fibre is regarded as a scheme over \(\mathbf F_{q^m}\), with the \(q^m\)-power Frobenius as its scheme Frobenius.
\end{enumerate}
\end{bpstatementit}

% proof: proof/lemma-comparison-via-relative-frobenius
\begin{proof}[{Proof}]\phantomsection\label{bp56}%
The only geometric fact we use is that, since \(r_m\) is a universal homeomorphism, pullback along it is an equivalence and its unit is an isomorphism. The rest is base change and composition of units, and each part uses the previous ones.

\emph{Proof of (1).} The identities \(p_mr_m=\Phi_X^m\) and \(f'_mr_m=f\) hold by the definition of \(r_m=(f,\Phi_X^m)\), since \(f\Phi_X^m=\Phi_Y^mf\). The map \(r_m\) is the geometric base change of the relative \(q^m\)-power Frobenius of \(X_0/Y_0\). It is integral, radicial and surjective [Stacks, \href{https://stacks.math.columbia.edu/tag/0CCB}{0CCB}], and finite because it is integral and of finite type. For affine \(Y_0=\operatorname{Spec}R\) and \(X_0=\operatorname{Spec}S\), with \(F\) the \(q\)-power Frobenius of \(R\), it is given on coordinate rings by

\[
R\otimes_{F^m,R}S\longrightarrow S,
\qquad b\otimes a\longmapsto f^\#(b)a^{q^m}.
\]

By topological invariance for classical objects of finite Tor amplitude, \(r_m^*\) is an equivalence, so the unit \(\eta_{r_m}\) in the construction is an isomorphism. In the construction, the first arrow is base change for compact supports, and the last comes from the composition of direct images with compact support, since \(r_m\) is proper. Moreover, pasting the cartesian squares defining \(X'_m\) and \(X'_n\), we get the comparison for \(m+n\), because the relative Frobenii compose and the unit of their composite is the composite of the units. Thus

\[
\theta_{f,m+n,K}=
\theta_{f,m,(\Phi_X^n)^*K}\circ
(\Phi_Y^m)^*\theta_{f,n,K},
\]

which is the compatibility with iteration.

\emph{Proof of (2).} The formula for \(\psi'_m\) follows by induction on \(m\) from the displayed iteration identity and the naturality of \(\theta\).

\emph{Proof of (3).} Both endomorphisms are a pullback map followed by \(R\Gamma_c\) of the correspondence, so we must see that the pullback along \(\Phi_Y^m\) on \(R\Gamma_c(Y,L)\) becomes the pullback along \(\Phi_X^m\) on \(R\Gamma_c(X,K)\). For a finite map \(h\) we also write \(h^*\) for the pullback map on cohomology with compact support, defined by the unit of \(h^*\dashv Rh_*\). Since base change is compatible with these units, the pullback along \(\Phi_Y^m\), followed by the base change isomorphism, is the pullback along \(p_m\), and the next arrow in \(\theta_{f,m}\) becomes \(r_m^*\). Their composite is \(r_m^*p_m^*=(p_mr_m)^*=(\Phi_X^m)^*\), and composing with \(R\Gamma_c(\psi_m)\) we see that the operators are intertwined.

\emph{Proof of (4).} We pull the construction back to the fixed point \(y\). Since \(\Phi_Y^m\) fixes \(\bar y\), the fibre of \(X'_m\to Y\) over \(\bar y\) is again the fibre \(X_{\bar y}\) of \(f\), and \(r_m\) restricts to a self-map of \(X_{\bar y}\). For \(f:\mathbf A^1\to s\) and \(m=1\), for instance, \(X'_1=\mathbf A^1\) and \(r_1\) is \(t\mapsto t^q\). In general, the base change of \(r_m\) to this fibre is its \(q^m\)-power scheme Frobenius, and by proper base change its unit restricts to the unit on the fibre. Hence the same calculation identifies \(\psi'_{m,\bar y}\) with the Frobenius operator on the cohomology with compact support of the fibre. Moreover, at a rational point \(x\) of the fibre, the restricted correspondence \(\psi_m|_{X_{\bar y}}\) has stalk map \(\psi_{m,\bar x}\), by composition of pullbacks.
\end{proof}

% text: prose-12-18
\phantomsection\label{bp57}%
To reduce the trace formula to pieces of \(X_0\), we cut \(X_0\) into an open subscheme and its closed complement.

% object: lemma-additivity-open-subscheme-complement (blueprint/chap-The-trace-formula/sec-Frobenius-and-Rf/lemma-additivity-open-subscheme-complement/_section.md)
\phantomsection\label{bp58}%
\begin{bpstatementit}{Lemma~12.19}{ (Additivity for an open subscheme and its complement)}
Let \(j_0:U_0\hookrightarrow X_0\) be open with reduced closed complement \(i_0:Z_0\to X_0\), and let \(K\in D^b_{\mathrm{ctf}}(X,A)\). A Frobenius correspondence \(\psi\) on \(K\) induces Frobenius correspondences on \(j^*K\), \(i^*K\), \(j_!j^*K\) and \(Ri_*i^*K\), whose stalk maps are the restrictions of \(\psi_m\). Writing \(T_{c_m}\) also for the endomorphisms of \(R\Gamma_c\) induced by the Frobenius correspondences on \(U\) and \(Z\), for every \(m\geq1\)

\[
\operatorname{tr}_A(T_{c_m}(\psi_m)\mid R\Gamma_c(X,K))
=\operatorname{tr}_A(T_{c_m}\mid R\Gamma_c(U,j^*K))
+\operatorname{tr}_A(T_{c_m}\mid R\Gamma_c(Z,i^*K)).
\]
\end{bpstatementit}

% proof: proof/lemma-additivity-open-subscheme-complement
\begin{proof}[{Proof}]\phantomsection\label{bp59}%
The four induced correspondences come from base change, and \(R\Gamma_c\) of the localization triangle is a two-step filtered object with an endomorphism, on which traces add.

Since Frobenius preserves opens defined over the finite field, its square with \(j\) is cartesian, and by base change we get the correspondences on \(j^*K\) and \(j_!j^*K\). The Frobenius square of the closed complement, however, is cartesian only after passing to reductions, and we obtain the correspondences on \(i^*K\) and \(Ri_*i^*K\) from Lemma lemma-comparison-via-relative-frobenius–(2) and nilpotent invariance. We then apply the functorial localization triangle \(j_!j^*K\to K\to Ri_*i^*K\) to the Frobenius correspondence and take direct images with compact support. Since the triangle is functorial, the Frobenius correspondence becomes an endomorphism of a two-step filtered object, and hence the traces add by Proposition \bpout{11.2}. Finally, Lemma lemma-comparison-via-relative-frobenius for \(j\) and \(i\) identifies the endomorphisms \(T_{c_m}\) of the cohomology of \(j_!j^*K\) and \(Ri_*i^*K\) with those on \(U\) and \(Z\).
\end{proof}


% object: sec-Finite-coefficients-4 (blueprint/chap-The-trace-formula/sec-Finite-coefficients-4/_section.md), content hidden
\subsection*{\texorpdfstring{12.4\quad Finite coefficients}{12.4 Finite coefficients}}
\phantomsection\label{bp60}%
\addcontentsline{toc}{subsection}{\texorpdfstring{12.4\quad Finite coefficients}{12.4 Finite coefficients}}
% text: prose-12-19
\phantomsection\label{bp61}%
The proof uses a single compactification, \(\mathbf A^1\subset\mathbf P^1\), so the Lefschetz–Verdier formula and the curve computation of Proposition \hyperref[bp49]{12.16} are applied only on \(\mathbf P^1\). This suffices because every affine scheme of finite type over \(\mathbf F_q\) is a closed subscheme of some \(\mathbf A^d\), and the fibres of a coordinate projection lie in \(\mathbf A^1\). On \(\mathbf P^1\) the only new fixed point is the rational point \(\infty\), where \(j_!K\) has zero stalk. The fibres over points of \(Y_0(\mathbf F_{q^m})\) are schemes over \(\mathbf F_{q^m}\), so we prove the theorem for all finite fields at once. The other coefficient rings are treated in Section \hyperref[bp64]{12.5}.

The base of the induction is the case of finite schemes. There we can see directly what the formula counts. For example, take \(X_0=\operatorname{Spec}\mathbf F_{q^2}\) and \(K=A\). Then the scheme \(X\) consists of two points interchanged by \(\Phi_X\). Hence \(T_{c_1}\) is the swap of \(A\oplus A\), with trace \(0=\#X_0(\mathbf F_q)\), and \(T_{c_2}\) is the identity, with trace \(2=\#X_0(\mathbf F_{q^2})\). In general \(T_{c_m}\) is a block permutation matrix, and only the points fixed by \(\Phi_X^m\) contribute to its trace.

The iterates \(\psi_m\), the endomorphisms \(T_{c_m}(\psi_m)\) and the stalk maps \(\psi_{m,\bar x}\) are those of Definition \hyperref[bp7]{12.2}.

% object: theorem-trace-formula-frobenius-correspondence (blueprint/chap-The-trace-formula/sec-Finite-coefficients-4/theorem-trace-formula-frobenius-correspondence/_section.md)
\phantomsection\label{bp62}%
\begin{bpstatementit}{Theorem~12.20}{ (Trace formula for a Frobenius correspondence)}
Let \(X_0/\mathbf F_q\) be separated and of finite type, put \(X=X_0\times_{\mathbf F_q}k\), let \(A\) be any finite commutative \(\ell\)-primary ring and \(K\in D^b_{\mathrm{ctf}}(X,A)\), and let \(\psi:\Phi_X^*K\to K\) be any Frobenius correspondence on \(K\). For every \(m\geq1\), in \(A\),

\[
\operatorname{tr}_A(T_{c_m}(\psi_m)\mid R\Gamma_c(X_{\mathrm{pro}\acute{e}t},K))
=\sum_{x\in X_0(\mathbf F_{q^m})}
\operatorname{tr}_A(\psi_{m,\bar x}\mid K_{\bar x}).
\]
\end{bpstatementit}

% proof: proof/theorem-trace-formula-frobenius-correspondence
\begin{proof}[{Proof}]\phantomsection\label{bp63}%
Both sides are defined, because the complexes are perfect by Lemma \bpout{10.1}. We prove the assertion simultaneously for all finite fields, all Frobenius correspondences and all \(m\), so that \(\mathbf F_q\) may be replaced by \(\mathbf F_{q^m}\) when we pass to a rational fibre. We go in four steps: finite schemes; the affine line, where we use the Lefschetz–Verdier formula on \(\mathbf P^1\); closed subschemes of \(\mathbf A^d\), by projecting away one coordinate; and arbitrary \(X_0\), by cutting it into affine pieces.

\emph{Finite schemes.} If \(X_0\) is finite, both sides of the formula can be read off from a block matrix. Indeed, up to nilpotents, \(X\) is the finite disjoint union of its geometric points. By nilpotent invariance, \(R\Gamma_c(X,K)=\bigoplus_{z\in X(k)}K_z\). In this decomposition, the pullback along \(\Phi_X^m\) given by the unit permutes these summands, and the block of \(T_{c_m}(\psi_m)\) from \(K_{\Phi_X^m(z)}\) to \(K_z\) is the stalk of \(\psi_m\) at \(z\), a map \(K_{\Phi_X^m(z)}=((\Phi_X^m)^*K)_z\to K_z\). On a finite projective model only the diagonal blocks contribute to the trace. These occur exactly at the points \(z\in X_0(\mathbf F_{q^m})\), with the required stalk maps.

\emph{The affine line.} This is the only step where we compute a local term. We pass to \(\mathbf P^1\), where the Lefschetz–Verdier formula applies, and extend \(K\) by zero, so that the new fixed point \(\infty\) contributes nothing. Let \(j:\mathbf A^1_k\hookrightarrow\mathbf P^1_k\). The correspondence \(\psi_m\) extends to a correspondence \(\bar\psi_m\) on \(j_!K\), compatible with iteration. Indeed, \(j\) is defined over \(\mathbf F_q\), so its square with the Frobenius maps is cartesian, and base change for this square gives \((\Phi_{\mathbf P^1}^m)^*j_!K\simeq j_!(\Phi_{\mathbf A^1}^m)^*K\). Lemma lemma-comparison-via-relative-frobenius, applied to \(j\), shows that the composition isomorphism

\[
R\Gamma_c(\mathbf A^1,K)
\simeq R\Gamma(\mathbf P^1,j_!K)
\]

intertwines the endomorphisms \(T_{c_m}\) induced by \(\psi\) and \(\bar\psi\). The graph is proper on the proper curve \(\mathbf P^1\), so by Theorem \bpout{11.13} the trace on the right is the sum of the local terms at the fixed points. By Proposition \hyperref[bp49]{12.16}, each local term is a trace on a stalk. In particular, the local term at infinity vanishes, because the stalk of \(j_!K\) there is zero. The other fixed points and stalk maps are those of \(\mathbf A^1_{\mathbf F_q}\). This proves the formula for the affine line.

\emph{Closed subschemes of the affine line.} Let \(X_0=\operatorname{Spec}\mathbf F_q[t]/I\). If \(I=0\), the previous paragraph applies. Otherwise \(I\) contains a nonzero polynomial, so \(X_0\) is finite over \(\mathbf F_q\) and the case of finite schemes applies. This also covers nonreduced subschemes.

\emph{Induction on affine coordinates.} We push \(K\) forward along a coordinate projection. This does not change the global side. The local side over a point of the base becomes the global side of a fibre lying in \(\mathbf A^1\). Suppose that the assertion holds for closed subschemes of \(\mathbf A^{d-1}\) over every finite field, and let \(X_0\hookrightarrow\mathbf A^d_{\mathbf F_q}\) be closed. Let \(f_0:X_0\to Y_0=\mathbf A^{d-1}_{\mathbf F_q}\) forget the last coordinate, and let \(f:X\to Y\) be its base change to \(k\). The complex \(L=Rf_!K\) lies in \(D^b_{\mathrm{ctf}}(Y,A)\) by Proposition \bpout{7.1}, and we equip it with the Frobenius correspondence \(\psi'\) of Lemma lemma-comparison-via-relative-frobenius. As Lemma lemma-comparison-via-relative-frobenius identifies the endomorphisms \(T_{c_m}\), the induction hypothesis shows that

\[
\operatorname{tr}_A(T_{c_m}(\psi_m)\mid R\Gamma_c(X,K))
=\sum_{y\in Y_0(\mathbf F_{q^m})}
\operatorname{tr}_A(\psi'_{m,\bar y}\mid L_{\bar y}).
\]

By Lemma lemma-comparison-via-relative-frobenius, each summand is the trace on the cohomology with compact support of the fibre \(X_{0,y}\), a closed subscheme of \(\mathbf A^1_{\mathbf F_{q^m}}\). Here the endomorphism is given by the restriction of \(\psi_m\) to the fibre and the first power of the Frobenius over \(\mathbf F_{q^m}\). The case of closed subschemes of the affine line then turns this summand into the sum of the traces on stalks over \(X_{0,y}(\mathbf F_{q^m})\). As \(y\) varies, these sets partition \(X_0(\mathbf F_{q^m})\), which proves the affine case.

\emph{General schemes.} Both sides are additive for a decomposition into an open subscheme and its closed complement: the global side by Lemma \hyperref[bp58]{12.19}, and the local side because the rational points split. It is therefore enough to cut \(X_0\) into affine pieces. Choose a finite affine open cover \(U_1,\ldots,U_N\) of \(X_0\), let \(W_i=\bigcup_{j\leq i}U_j\) and \(V_i=W_i\setminus W_{i-1}\), with the reduced closed structure. Each \(V_i\) is closed in the affine \(U_i\), hence affine, and it is the closed complement of \(W_{i-1}\) in \(W_i\). Applying Lemma \hyperref[bp58]{12.19} successively, we see that the traces add over this stratification. The affine case applied to each \(V_i\) now gives the equality for all separated \(X_0\) of finite type.
\end{proof}


% object: sec-ell-adic-coefficients-5 (blueprint/chap-The-trace-formula/sec-ell-adic-coefficients-5/_section.md), content hidden
\subsection*{\texorpdfstring{12.5\quad $\ell$-adic coefficients}{12.5 ell-adic coefficients}}
\phantomsection\label{bp64}%
\addcontentsline{toc}{subsection}{\texorpdfstring{12.5\quad $\ell$-adic coefficients}{12.5 ell-adic coefficients}}
% text: prose-12-20
\phantomsection\label{bp65}%
To pass from finite coefficient rings to \(\mathcal O_E\), \(E\) and \(\overline{\mathbf Q}_\ell\) we follow the steps of Proposition \bpout{7.12}. Since the statement is now an equality of traces in \(\Lambda\) rather than an isomorphism, we carry the steps out directly. First, for \(\mathcal O_E\) we reduce modulo \(\pi_E^n\). This works because both sides are traces of perfect complexes, traces commute with reduction, and \(\mathcal O_E\) embeds in \(\varprojlim_n\mathcal O_E/\pi_E^n\). Next, for \(E\) we clear denominators: a Frobenius correspondence \(\psi_E\) with coefficients \(E\) need not come from an integral one; only some multiple \(\pi_E^N\psi_E\) does, and this integral model need not be invertible even when \(\psi_E\) is (Lemma \hyperref[bp69]{12.22}). This is why we prove the formula for correspondences that are not invertible. Finally, for \(\overline{\mathbf Q}_\ell\) we descend to a finite extension of \(\mathbf Q_\ell\).

The Frobenius comparison \(\kappa_m\) of Lemma \hyperref[bp9]{12.3} extends to the coefficient rings \(\mathcal O_E\), \(E\) and \(\overline{\mathbf Q}_\ell\). Let \(K\in\Dbc(X,\mathcal O_E)\) have tower of reductions \((K_n)\) modulo \(\pi_E^n\) (Definition \bpout{2.61}). Both \((\Phi_X^m)^*\) and \(\operatorname{Fr}_X^{m*}\) preserve derived completeness and commute with reduction, so each is the derived limit of its values on the \(K_n\). The finite comparisons \(\kappa_m\) come from a single functorial identification on the classical site. Hence they are compatible with reduction of coefficients, and their derived limit is the integral \(\kappa_m\). The finite comparisons are compatible with iteration and with the units of the pullbacks, and therefore so is the integral one. It follows that the integral \(\kappa_m\) identifies the endomorphisms \(T_{c_m}\) with the geometric Frobenius operators as in Lemma \hyperref[bp9]{12.3}. For \(E\) we localize on integral models. By the localized Hom comparison (9.21) the result does not depend on the choice of model. For \(\overline{\mathbf Q}_\ell\) we descend from finite extensions.

% object: proposition-integral-trace-formula (blueprint/chap-The-trace-formula/sec-ell-adic-coefficients-5/proposition-integral-trace-formula/_section.md)
\phantomsection\label{bp66}%
\begin{bpstatementit}{Proposition~12.21}{ (Integral trace formula)}
Let \(K\in\Dbc(X,\mathcal O_E)\) and let \(\theta:\Phi_X^*K\to K\) be any morphism, that is, a Frobenius correspondence with integral coefficients. Define its iterates \(\theta_m\), the endomorphisms \(T_{c_m}(\theta_m)\) and the stalk maps \(\theta_{m,\bar y}\) by the formulas of Definition \hyperref[bp7]{12.2}. Then for every \(m\geq1\) we have

\[
\operatorname{tr}_{\mathcal O_E}
\bigl(T_{c_m}(\theta_m)\mid R\Gamma_c(X_{\mathrm{pro}\acute{e}t},K)\bigr)
=\sum_{y\in X_0(\mathbf F_{q^m})}
\operatorname{tr}_{\mathcal O_E}
\bigl(\theta_{m,\bar y}\mid K_{\bar y}\bigr),
\tag{12.1}
\]

where both sides are categorical traces of perfect complexes over \(\mathcal O_E\). Their reductions modulo \(\pi_E^n\) are the corresponding traces of Theorem \hyperref[bp62]{12.20} for \(K_n\) and the reduction \(\bar\theta_n:\Phi_X^*K_n\to K_n\) of \(\theta\).
\end{bpstatementit}

% proof: proof/proposition-integral-trace-formula
\begin{proof}[{Proof}]\phantomsection\label{bp67}%
We show that both sides reduce modulo \(\pi_E^n\) to the two sides of the finite formula, by identifying the reductions of the complexes, of the operators and of the traces. This suffices, since an element of \(\mathcal O_E\) that vanishes modulo every \(\pi_E^n\) is zero.

\emph{Complexes and operators.} By Lemma \bpout{10.1}, \(R\Gamma_c(X_{\mathrm{pro}\acute{e}t},K)\) and the geometric stalks are perfect, and Proposition \bpout{7.10} provides canonical isomorphisms

\[
\begin{aligned}
R\Gamma_c(X_{\mathrm{pro}\acute{e}t},K)\otimes^{\mathbf L}_{\mathcal O_E}\mathcal O_E/\pi_E^n
&\simeq R\Gamma_c(X_{\mathrm{pro}\acute{e}t},K_n),\\
K_{\bar y}\otimes^{\mathbf L}_{\mathcal O_E}\mathcal O_E/\pi_E^n&\simeq(K_n)_{\bar y}.
\end{aligned}
\]

Under these isomorphisms the reductions of the integral operators are the operators of \(\bar\theta_n\), because extension of scalars commutes with the pullbacks, with the units defining the endomorphisms \(T_{c_m}\), and with composition, and hence with the iterates \(\theta_m\). Moreover Theorem \hyperref[bp62]{12.20} applies for every \(n\), because \(K_n\in D^b_{\mathrm{ctf}}(X,\mathcal O_E/\pi_E^n)\).

\emph{Traces.} Represent each of the two integral complexes by a bounded complex \(P^\bullet\) of finite projective \(\mathcal O_E\)-modules and its endomorphism by a chain map \(v^\bullet\). By Proposition \bpout{11.2} its trace is

\[
\operatorname{tr}_{\mathcal O_E}(v)
=\sum_i(-1)^i\operatorname{Tr}(v^i\mid P^i).
\]

The reduction of this expression modulo \(\pi_E^n\) is the trace on \(P^\bullet\otimes_{\mathcal O_E}\mathcal O_E/\pi_E^n\), as each \(P^i\) is finite projective.

\emph{Conclusion.} The reduction of each side of (12.1) is therefore the corresponding side of the finite formula for \((K_n,\bar\theta_n)\). Hence the difference of the two sides lies in \(\pi_E^n\mathcal O_E\) for every \(n\). It is zero because \(\mathcal O_E\hookrightarrow\varprojlim_n\mathcal O_E/\pi_E^n\) is injective.
\end{proof}

% text: prose-12-21
\phantomsection\label{bp68}%
To pass from \(\mathcal O_E\) to \(E\) we need an integral Frobenius correspondence whose rationalization is the given one, up to a power of \(\pi_E\). The next lemma provides it and records how the operators of both sides of the trace formula change under this rescaling.

% object: lemma-integral-models-rational-frobenius-correspondences (blueprint/chap-The-trace-formula/sec-ell-adic-coefficients-5/lemma-integral-models-rational-frobenius-correspondences/_section.md)
\phantomsection\label{bp69}%
\begin{bpstatementit}{Lemma~12.22}{ (Integral models for rational Frobenius correspondences)}
Let \(K_E\in\Dbc(X,E)\) and let \(\psi_E:\Phi_X^*K_E\to K_E\) be any morphism. There exist \(K_0\in\Dbc(X,\mathcal O_E)\), an isomorphism \(K_0\otimes^{\mathbf L}_{\mathcal O_E}E\simeq K_E\), an integer \(N\geq0\), and a Frobenius correspondence \(\theta:\Phi_X^*K_0\to K_0\) such that

\[
\psi_E=\pi_E^{-N}(\theta\otimes^{\mathbf L}_{\mathcal O_E}E).
\]

For every \(m\geq1\) the operators of Definition \hyperref[bp7]{12.2} satisfy

\[
\begin{aligned}
T_{c_m}((\psi_E)_m)&=\pi_E^{-Nm}
\bigl(T_{c_m}(\theta_m)\otimes^{\mathbf L}_{\mathcal O_E}E\bigr),\\
(\psi_E)_{m,\bar y}&=\pi_E^{-Nm}
\bigl(\theta_{m,\bar y}\otimes^{\mathbf L}_{\mathcal O_E}E\bigr).
\end{aligned}
\tag{12.2}
\]
\end{bpstatementit}

% proof: proof/lemma-integral-models-rational-frobenius-correspondences
\begin{proof}[{Proof}]\phantomsection\label{bp70}%
Morphisms in the rational category are integral morphisms divided by a power of \(\pi_E\). This gives \(K_0\) and \(\theta\). The formulas for the operators hold because the operators of Definition \hyperref[bp7]{12.2} commute with extension of scalars.

By Proposition \bpout{3.27}, extension of scalars induces an equivalence

\[
\Dbc(X,\mathcal O_E)[1/\ell]\simeq\Dbc(X,E),
\]

where the source has the objects of \(\Dbc(X,\mathcal O_E)\) and Hom groups tensored with \(E\). By essential surjectivity, \(K_E\) has an integral model \(K_0\). The model constructed there is built by finite cones and shifts from the integral generators of Lemma lemma-integral-models, and hence in general it is not a lattice in a lisse sheaf.

Since inverse image preserves integral constructibility, \(\Phi_X^*K_0\) is an integral model of \(\Phi_X^*K_E\). By the description of the Hom groups above, applied to the pair \((\Phi_X^*K_0,K_0)\), we may write \(\psi_E=\pi_E^{-N}(\theta\otimes^{\mathbf L}_{\mathcal O_E}E)\). Since pullback is \(\mathcal O_E\)-linear, each of the \(m\) factors of the iterate contributes \(\pi_E^{-N}\), and so \((\psi_E)_m=\pi_E^{-Nm}(\theta_m\otimes^{\mathbf L} E)\). By Proposition \bpout{7.10}, the change-of-coefficients maps for cohomology with compact support and for stalks respect the units of the pullback adjunctions. Hence the pullback map \((\Phi_X^m)^*\) on \(R\Gamma_c\) for \(E\) is the extension of scalars of the integral one. The same holds for the identification \(((\Phi_X^m)^*K_E)_{\bar y}\simeq (K_E)_{\bar y}\) at a fixed point \(\bar y\) of \(\Phi_X^m\), because the comparison for inverse images is compatible with composition and with extension of scalars. Applying \(T_{c_m}\) and taking stalks at \(\bar y\) in the formula for \((\psi_E)_m\) gives (12.2).
\end{proof}

% object: remark-integral-frobenius-correspondences-not-invertible (blueprint/chap-The-trace-formula/sec-ell-adic-coefficients-5/remark-integral-frobenius-correspondences-not-invertible/_section.md)
\phantomsection\label{bp71}%
\begin{bpstatement}{Remark~12.23}{ (Integral Frobenius correspondences need not be invertible)}
In Lemma \hyperref[bp69]{12.22}, \(K_0\) need not be a lattice in a lisse sheaf, and \(\theta\) need not be invertible even when \(\psi_E\) is. The factor \(\pi_E^{-N}\) is needed because \(\psi_E\) itself need not be integral: on a point, the automorphism \(\ell^{-1}\) of \(E\) preserves no lattice. Even with this factor, \(\theta\) cannot always be chosen invertible, already on a point. On \(E^2\), the two eigenvalues of \(\operatorname{diag}(1,\ell)\) cannot both become units after multiplying the operator by a single scalar. So Lemma \hyperref[bp69]{12.22} only makes \(\pi_E^N\psi_E\) integral, and the reductions of \(\theta\) may fail to be invertible. For this reason Theorem \hyperref[bp62]{12.20} is stated for arbitrary Frobenius correspondences.
\end{bpstatement}

% text: prose-12-22
\phantomsection\label{bp72}%
With the integral formula and integral models of rational Frobenius correspondences, we obtain the trace formula stated at the beginning of this chapter for all coefficient rings.

% object: theorem-grothendieck-lefschetz-trace-formula (blueprint/chap-The-trace-formula/sec-ell-adic-coefficients-5/theorem-grothendieck-lefschetz-trace-formula/_section.md)
\phantomsection\label{bp73}%
\begin{bpstatementit}{Theorem~12.24}{ (Grothendieck–Lefschetz trace formula)}
Let \(X_0\) be separated and of finite type over \(\mathbf F_q\), put \(X=X_0\times_{\mathbf F_q}\overline{\mathbf F}_q\), and suppose that \(\ell\nmid q\). Let

\[
\Lambda\in\{A,\mathbf Z_\ell,\mathcal O_E,\mathbf Q_\ell,E,\overline{\mathbf Q}_\ell\},
\qquad K\in D^b_{\mathrm{ctf}}(X,\Lambda),
\]

with a Weil structure as in Definition \bpout{10.18}. Here \(A\) is any finite commutative \(\ell\)-primary ring, and for integral and rational coefficients, \(D^b_{\mathrm{ctf}}(X,\Lambda)=\Dbc(X,\Lambda)\). Let \(\operatorname{Frob}_{X,K}\) be the geometric Frobenius on \(R\Gamma_c(X_{\mathrm{pro}\acute{e}t},K)\), and let \(\operatorname{Frob}_{y,m}\) be the Frobenius acting on the stalk \(K_{\bar y}\) at \(y\in X_0(\mathbf F_{q^m})\). Then for every \(m\geq1\) we have, in \(\Lambda\),

\[
\operatorname{tr}_{\Lambda}
\bigl(\operatorname{Frob}_{X,K}^m\mid R\Gamma_c(X_{\mathrm{pro}\acute{e}t},K)\bigr)
=\sum_{y\in X_0(\mathbf F_{q^m})}
\operatorname{tr}_{\Lambda}
\bigl(\operatorname{Frob}_{y,m}\mid K_{\bar y}\bigr).
\tag{12.3}
\]

The complexes inside the traces are perfect over \(\Lambda\), and the formula is compatible with the change-of-coefficients maps of Proposition \bpout{7.10}. For \(\Lambda=\mathbf Q_\ell,E,\overline{\mathbf Q}_\ell\), it reads

\[
\sum_i(-1)^i\operatorname{Tr}
\bigl(\operatorname{Frob}_{X,K}^m\mid H_c^i(X_{\mathrm{pro}\acute{e}t},K)\bigr)
=\sum_{y\in X_0(\mathbf F_{q^m})}\sum_j(-1)^j
\operatorname{Tr}\bigl(\operatorname{Frob}_{y,m}\mid H^j(K_{\bar y})\bigr),
\tag{12.4}
\]

where all sums over degrees are finite.
\end{bpstatementit}

% text: prose-12-23
\phantomsection\label{bp74}%
For finite and integral rings, (12.3) is an identity of categorical traces of perfect complexes, which makes sense even when the cohomology modules are not projective.

% proof: proof/theorem-grothendieck-lefschetz-trace-formula
\begin{proof}[{Proof of Theorem 12.24}]\phantomsection\label{bp75}%
In each case we turn the Weil structure into a Frobenius correspondence by composing with \(\kappa_1\), and we use Lemma \hyperref[bp9]{12.3} to see that this does not change the operators. We then treat the coefficient rings in the order finite, integral, rational, algebraic, each case using the previous one.

For finite \(A\), we apply Theorem \hyperref[bp62]{12.20} to \(\psi=\phi\kappa_1\). Then Lemma \hyperref[bp9]{12.3} identifies \(T_{c_m}(\psi_m)\) with \(\operatorname{Frob}_{X,K}^m\) and, at each fixed point, \(\psi_{m,\bar y}\) with \(\operatorname{Frob}_{y,m}\). Similarly, for \(\mathcal O_E\), composing the Weil structure with \(\kappa_1\) we get a Frobenius correspondence \(\Phi_X^*K\to K\). By the integral form of Lemma \hyperref[bp9]{12.3}, the endomorphisms it induces on \(R\Gamma_c\) and on the stalks are those induced by the Weil structure. The identity therefore follows from Proposition \hyperref[bp66]{12.21}. In particular, taking \(E=\mathbf Q_\ell\) gives the case of \(\mathbf Z_\ell\).

For coefficients \(E\), the Weil structure composed with \(\kappa_1\) is a Frobenius correspondence \(\psi_E\), and by the rational form of Lemma \hyperref[bp9]{12.3} the endomorphisms it induces on \(R\Gamma_c\) and on the stalks are those induced by the Weil structure. We apply Lemma \hyperref[bp69]{12.22} to \(\psi_E\), apply (12.1) to \((K_0,\theta)\), extend scalars to \(E\), and multiply by \(\pi_E^{-Nm}\). By linearity of categorical traces and their compatibility with extension of scalars (Proposition \bpout{11.2}), together with (12.2), this gives the identity for every \(m\). The case \(E=\mathbf Q_\ell\) gives the case of \(\mathbf Q_\ell\).

For \(\overline{\mathbf Q}_\ell\), complexes and morphisms are colimits of their versions over finite extensions, so it is enough to find the Weil complex over some finite \(E\) and to check that the answer does not depend on the choice. By Proposition \bpout{7.11} the complex \(K\), its Weil structure \(\phi\) and \(\phi^{-1}\) descend to a common finite extension \(E/\mathbf Q_\ell\). After a further finite extension the two identities saying that these maps are inverse to each other hold as well. Thus we get a Weil complex at a finite stage, and by extension of scalars its formula over \(E\) gives (12.3). If we choose another representative, any isomorphism compatible with the Weil structures descends to a common finite extension, so the resulting operators are conjugate. Since the trace is compatible with extension of scalars and invariant under conjugation, the result does not depend on the stage or on the representative.

The complexes are perfect by Lemma \bpout{10.1}. Compatibility with the change-of-coefficients maps holds because the pullbacks, the endomorphisms \(T_{c_m}\) and the iterated Frobenius correspondences commute with the comparisons of Proposition \bpout{7.10}. Finally, over a field the categorical trace is the alternating trace on cohomology by Proposition \bpout{11.2}, and this proves (12.4).
\end{proof}


% object: sec-The-function-sheaf-dictionary-2 (blueprint/chap-The-trace-formula/sec-The-function-sheaf-dictionary-2/_section.md), content hidden
\subsection*{\texorpdfstring{12.6\quad The function–sheaf dictionary}{12.6 The function–sheaf dictionary}}
\phantomsection\label{bp76}%
\addcontentsline{toc}{subsection}{\texorpdfstring{12.6\quad The function–sheaf dictionary}{12.6 The function–sheaf dictionary}}
% text: prose-12-24
\phantomsection\label{bp77}%
Under the trace formula, \(Rf_!\) becomes summation over fibres and \(\otimes\) becomes multiplication of functions on rational points. This is Grothendieck's function–sheaf dictionary [\hyperref[bp106]{SGA4½}, Sommes trig.] (Section \bpout{1.3}), and in Part III we compute with sheaves through their functions in this way.

% object: definition-trace-function (blueprint/chap-The-trace-formula/sec-The-function-sheaf-dictionary-2/definition-trace-function/_section.md)
\phantomsection\label{bp78}%
\begin{bpstatement}{Definition~12.25}{ (Trace function)}
Let \((K,\phi)\) be a Weil complex on \(X\) as in Theorem \hyperref[bp73]{12.24}, and let \(m\geq1\). The \emph{trace function} of \(K\) is

\[
t_{K,m}:X_0(\mathbf F_{q^m})\longrightarrow\Lambda,\qquad
t_{K,m}(x)=\operatorname{tr}_\Lambda\bigl(\operatorname{Frob}_{x,m}
\mid K_{\bar x}\bigr).
\]
\end{bpstatement}

% text: prose-12-25
\phantomsection\label{bp79}%
For \(\Lambda\) a field, \(t_{K,m}(x)\) is the alternating sum of the traces of Frobenius on the stalks \(\mathcal H^i(K)_{\bar x}\) (Proposition \bpout{11.2}). For example, the constant sheaf \(\Lambda_X\) with its constant Weil structure has trace function \(1\). In terms of trace functions, the theorem says that \(\operatorname{tr}_\Lambda(\operatorname{Frob}_{X,K}^m
\mid R\Gamma_c(X_{\mathrm{pro}\acute{e}t},K))\) is the sum of \(t_{K,m}\) over \(X_0(\mathbf F_{q^m})\). Applied to the fibres of a morphism, it shows that direct image with compact support sums trace functions over the fibres.

% object: corollary-trace-functions-direct-image-compact (blueprint/chap-The-trace-formula/sec-The-function-sheaf-dictionary-2/corollary-trace-functions-direct-image-compact/_section.md)
\phantomsection\label{bp80}%
\begin{bpstatementit}{Corollary~12.26}{ (Trace functions and direct image with compact support)}
Let \(f_0:X_0\to Y_0\) be a morphism of separated \(\mathbf F_q\)-schemes of finite type, with geometric base change \(f:X\to Y\), and let \((K,\phi)\) be a Weil complex on \(X\) with \(K\in D^b_{\mathrm{ctf}}(X,\Lambda)\) for an admissible \(\Lambda\). Give \(Rf_!K\in D^b_{\mathrm{ctf}}(Y,\Lambda)\) the Weil structure induced by \(\phi\) (after Definition \bpout{10.18}). Then for every \(m\geq1\) and \(y\in Y_0(\mathbf F_{q^m})\),

\[
t_{Rf_!K,m}(y)=\sum_{\substack{x\in X_0(\mathbf F_{q^m})\\ f_0(x)=y}}
t_{K,m}(x).
\]
\end{bpstatementit}

% proof: proof/corollary-trace-functions-direct-image-compact
\begin{proof}[{Proof}]\phantomsection\label{bp81}%
The left side is a trace on a stalk of \(Rf_!K\). By base change that stalk is the compactly supported cohomology of a fibre. Hence the corollary is the trace formula for the fibre, once we check that the Frobenius operators correspond.

Fix \(y\) and let \(X_{0,y}\) be the fibre of \(f_0\) over \(y\). It is a separated scheme of finite type over \(\mathbf F_{q^m}\), and its rational points are the \(x\in X_0(\mathbf F_{q^m})\) with \(f_0(x)=y\). Restrict \(K\) to its geometric fibre \(X_{\bar y}\) and give the restriction the Weil structure over \(\mathbf F_{q^m}\) obtained from the \(m\)-th iterate of \(\phi\). By base change for compact supports, \((Rf_!K)_{\bar y}\) is \(R\Gamma_c(X_{\bar y},K|_{X_{\bar y}})\), and \(\operatorname{Frob}_{y,m}\) is the geometric Frobenius operator of this Weil complex. For finite \(A\) this is part (4) of Lemma \hyperref[bp55]{12.18}, applied to the Frobenius correspondence \(\phi\kappa_1\) and translated by Lemma \hyperref[bp9]{12.3}. By the comparisons of Proposition \bpout{7.10} it extends to the other coefficient rings. Since the Frobenius acting on the stalk of the Weil complex on the fibre at a rational point \(x\) of \(X_{0,y}\) is \(\operatorname{Frob}_{x,m}\), the stated identity is Theorem \hyperref[bp73]{12.24} for \(X_{0,y}\) over \(\mathbf F_{q^m}\), with exponent \(1\).
\end{proof}

% object: example-trace-function-squaring (blueprint/chap-The-trace-formula/sec-The-function-sheaf-dictionary-2/example-trace-function-squaring/_section.md)
\phantomsection\label{bp82}%
\begin{bpstatement}{Example~12.27}{ (The squaring map)}
Let \(q\) be odd, let \(f_0:\mathbf A^1_{\mathbf F_q}\to\mathbf A^1_{\mathbf F_q}\) be \(x\mapsto x^2\), and let \(K=\Lambda_X\) with its constant Weil structure, so that \(t_{K,m}=1\). The fibre of \(f_0\) over \(y\in\mathbf F_{q^m}\) has \(1+\chi_m(y)\) rational points, where \(\chi_m\) is the quadratic character of \(\mathbf F_{q^m}^*\), extended by \(\chi_m(0)=0\). Hence

\[
t_{Rf_!\Lambda,m}(y)=1+\chi_m(y).
\]

The trace function of \(Rf_!\Lambda\) is not constant, although the sheaf we pushed forward is: over \(\mathbf G_m\) the map is a finite étale double cover, and \(\chi_m(y)\) records whether the fibre over \(y\) splits into two rational points or is a single point of degree \(2\). Summing over \(y\) recovers the trace formula on the source. Since \(\chi_m\) takes the values \(\pm1\) equally often on \(\mathbf F_{q^m}^*\), the sum is \(q^m=\#\mathbf A^1(\mathbf F_{q^m})\).
\end{bpstatement}

% text: prose-12-26
\phantomsection\label{bp83}%
The other operations on functions also have counterparts on sheaves. These counterparts, together with the corollary, are the entries of the dictionary that we use. In this dictionary, morphisms come from morphisms of separated \(\mathbf F_q\)-schemes of finite type, and the operations carry the Weil structures described after Definition \bpout{10.18}.

% object: proposition-function-sheaf-dictionary (blueprint/chap-The-trace-formula/sec-The-function-sheaf-dictionary-2/proposition-function-sheaf-dictionary/_section.md)
\phantomsection\label{bp84}%
\begin{bpstatementit}{Proposition~12.28}{ (The function–sheaf dictionary)}
Let \(\Lambda\) be as in Theorem \hyperref[bp73]{12.24}, let \(m\geq1\), and let \(K\), \(L\) be Weil complexes in \(D^b_{\mathrm{ctf}}(X,\Lambda)\).

\begin{enumerate}
\item For \(g_0:W_0\to X_0\) and \(w\in W_0(\mathbf F_{q^m})\), \(t_{g^*K,m}(w)=t_{K,m}(g_0(w))\).
\item \(t_{K\otimes^{\mathbf L}_\Lambda L,m}=t_{K,m}\,t_{L,m}\).
\item For \(f_0:X_0\to Y_0\) and \(y\in Y_0(\mathbf F_{q^m})\), \(t_{Rf_!K,m}(y)=\sum_{x\in X_0(\mathbf F_{q^m}),\,f_0(x)=y}t_{K,m}(x)\).
\item For \(U_0\subset X_0\) open with closed complement \(Z_0\), and \(j:U\to X\), \(i:Z\to X\) the inclusions, \(t_{K,m}=t_{j_!j^*K,m}+t_{i_*i^*K,m}\), where \(t_{j_!j^*K,m}\) is \(t_{K,m}\) on \(U_0(\mathbf F_{q^m})\) and \(0\) on \(Z_0(\mathbf F_{q^m})\).
\item \(\sum_{x\in X_0(\mathbf F_{q^m})}t_{K,m}(x)
  =\operatorname{tr}_\Lambda\bigl(\operatorname{Frob}_{X,K}^m\mid
  R\Gamma_c(X_{\mathrm{pro}\acute{e}t},K)\bigr)\).
\end{enumerate}
\end{bpstatementit}

% proof: proof/proposition-function-sheaf-dictionary
\begin{proof}[{Proof}]\phantomsection\label{bp85}%
Part (5) is Theorem \hyperref[bp73]{12.24} and part (3) is Corollary \hyperref[bp80]{12.26}.

For (1), the stalk of \(g^*K\) at \(\bar w\) is \(K_{\bar x}\) with \(x=g_0(w)\), and the Weil structure of \(g^*K\) is \(g^*\phi\) composed with the base change isomorphism. Under this identification of stalks the \(m\)-th iterate of the Weil structure of \(g^*K\) at \(\bar w\) is the \(m\)-th iterate of \(\phi\) at \(\bar x\), so the two Frobenius endomorphisms agree.

For (2), the stalk of \(K\otimes^{\mathbf L}_\Lambda L\) at \(\bar x\) is \(K_{\bar x}\otimes^{\mathbf L}_\Lambda L_{\bar x}\), since pullback to \(\bar x\) commutes with derived tensor products, and its Frobenius is the tensor product of the two Frobenius endomorphisms. These stalks are perfect (Lemma \bpout{10.1}), so we may represent them by bounded complexes \(P\), \(Q\) of finite projective modules and the endomorphisms by chain maps \(v\), \(w\). Now the term of degree \(n\) of \(P\otimes Q\) is \(\bigoplus_{i+j=n}P^i\otimes Q^j\), the trace of \(v^i\otimes w^j\) on \(P^i\otimes Q^j\) is \(\operatorname{Tr}(v^i)\operatorname{Tr}(w^j)\), and \((-1)^n=(-1)^i(-1)^j\). By the chain formula of Proposition \bpout{11.2} the trace on \(P\otimes Q\) is therefore the product of the traces on \(P\) and on \(Q\).

For (4), note that \(j_!=Rj_!\) and \(i_*=Ri_!\), with the same Weil structures. By (3), \(t_{j_!j^*K,m}(x)\) is the sum of \(t_{j^*K,m}\) over the fibre of \(j_0\) at \(x\). This fibre is the point \(x\) if \(x\in U_0\), and it is empty otherwise. By (1), \(t_{j^*K,m}(x)=t_{K,m}(x)\) for \(x\in U_0(\mathbf F_{q^m})\). The same argument applies to \(i_*i^*K\), and the two functions add up to \(t_{K,m}\).
\end{proof}

% text: prose-12-27
\phantomsection\label{bp86}%
The Fourier transform of Chapter \bpout{16} is one more entry, and the proposition shows why it corresponds to the discrete transform. Let \(p_1,p_2\) be the projections of \(\mathbf A^1\times\mathbf A^1\), let \(\mu(x,y)=xy\), let \(\psi:\mathbf F_q\to E^\times\) be a nontrivial additive character, and let \(\mathcal K_\psi=\mu^*\mathcal L_\psi\), where \(\mathcal L_\psi\) is the Artin–Schreier sheaf of Lemma \bpout{16.1} (proved in Chapter \bpout{16} below). Here \(\psi\) is a character, not a Frobenius correspondence. With this kernel the transform is \(\operatorname{FT}_\psi K=Rp_{2!}(p_1^*K\otimes\mathcal K_\psi)[1]\). For \(y\in\mathbf F_{q^m}\) we compute

\[
\begin{aligned}
t_{\operatorname{FT}_\psi K,m}(y)
&=-\,t_{Rp_{2!}(p_1^*K\otimes\mathcal K_\psi),m}(y)
&&\text{(shift)}\\
&=-\sum_{x\in\mathbf F_{q^m}}t_{p_1^*K\otimes\mathcal K_\psi,m}(x,y)
&&\text{by (3)}\\
&=-\sum_{x\in\mathbf F_{q^m}}t_{p_1^*K,m}(x,y)\,t_{\mu^*\mathcal L_\psi,m}(x,y)
&&\text{by (2)}\\
&=-\sum_{x\in\mathbf F_{q^m}}t_{K,m}(x)\,
\psi\bigl(\operatorname{Tr}_{\mathbf F_{q^m}/\mathbf F_q}(xy)\bigr)
&&\text{by (1) and Lemma~16.1.1.}
\end{aligned}
\]

The first equality holds because a shift changes the sign of a trace, by the chain formula of Proposition \bpout{11.2}. Each of the other three steps is one operation in the formula \(\hat t(y)=\sum_xt(x)\psi(xy)\): pull \(t\) back to the plane, multiply by the kernel, and sum over the fibres of \(p_2\).


% object: sec-Rationality-and-the-functional-equation (blueprint/chap-The-trace-formula/sec-Rationality-and-the-functional-equation/_section.md), content hidden
\subsection*{\texorpdfstring{12.7\quad Rationality and the functional equation}{12.7 Rationality and the functional equation}}
\phantomsection\label{bp87}%
\addcontentsline{toc}{subsection}{\texorpdfstring{12.7\quad Rationality and the functional equation}{12.7 Rationality and the functional equation}}
% text: prose-12-28
\phantomsection\label{bp88}%
Over a field of characteristic zero, the traces of all powers of an endomorphism determine its characteristic polynomial. Taking the trace formula for all \(m\geq1\), we therefore obtain an identity of power series: a product over closed points equals an alternating product of characteristic polynomials of Frobenius on cohomology. This step divides by \(m\), so in the determinant formulas below \(\Lambda\) is one of the fields \(\mathbf Q_\ell\), \(E\), \(\overline{\mathbf Q}_\ell\). Characteristic zero is needed: on \(\mathbf F_p^{\oplus p}\) every power of the identity has trace \(p=0\), as for the zero endomorphism, but \(\det(1-T\cdot\operatorname{id})=(1-T)^p=1-T^p\) is not \(1\).

For \(K=\Lambda_X\) the product over cohomology is finite, so \(Z(X_0,T)\) is rational, and for smooth proper \(X_0\) duality pairs the factors in degrees \(i\) and \(2d-i\) (Corollary \hyperref[bp100]{12.33}).

% object: lemma-logarithm-determinant (blueprint/chap-The-trace-formula/sec-Rationality-and-the-functional-equation/lemma-logarithm-determinant/_section.md)
\phantomsection\label{bp89}%
\begin{bpstatementit}{Lemma~12.29}{ (Logarithm of a determinant)}
Let \(L\) be a field of characteristic zero. For an endomorphism \(v\) of a finite-dimensional \(L\)-vector space \(V\) we have

\[
\det(1-Tv\mid V)^{-1}
=\exp\left(\sum_{m\geq1}
\operatorname{Tr}(v^m\mid V)\frac{T^m}{m}\right)
\quad\text{in }1+TL[[T]].
\tag{12.5}
\]

For an endomorphism \(v\) of a perfect complex \(P\) in \(D(L)\), put

\[
\det(1-Tv\mid P)
:=\prod_i\det(1-Tv\mid H^i(P))^{(-1)^i}.
\]

Then (12.5) holds for \(P\), with \(\operatorname{Tr}(v^m\mid V)\) replaced by the categorical trace \(\operatorname{tr}_L(v^m\mid P)\).
\end{bpstatementit}

% proof: proof/lemma-logarithm-determinant
\begin{proof}[{Proof}]\phantomsection\label{bp90}%
For a single eigenvalue \(\alpha\) the identity is \(-\log(1-\alpha T)=\sum_m\alpha^mT^m/m\), and the general case could be reduced to this by triangularizing \(v\) over an algebraic closure. Instead we use logarithmic derivatives, which allows us to stay over \(L\). In \(L[[T]]\), the matrix \(1-Tv\) is invertible, with inverse \(\sum_{r\geq0}T^rv^r\). Moreover, the logarithmic derivative of the determinant of an invertible matrix \(M\) over \(L[[T]]\) is \(\operatorname{Tr}(M^{-1}M')\), and hence

\[
\frac{d}{dT}\log\det(1-Tv)
=-\operatorname{Tr}\bigl(v(1-Tv)^{-1}\bigr)
=-\sum_{m\geq1}\operatorname{Tr}(v^m)T^{m-1}.
\]

Since both logarithms have constant coefficient zero and \(L\) has characteristic zero, we may integrate and exponentiate, and we obtain (12.5). For a perfect complex, the categorical trace over a field is the alternating sum of the traces on the finitely many nonzero cohomology spaces (Proposition \bpout{11.2}), and the assertion follows by multiplying the identities for the \(H^i(P)\).
\end{proof}

% text: prose-12-29
\phantomsection\label{bp91}%
The local factors of the Euler product are indexed by closed points. A closed point \(x\in|X_0|\) of degree \(d\) over \(\mathbf F_q\) corresponds to a Frobenius orbit of \(d\) geometric points, and its Frobenius on the stalk is the Frobenius going once around this orbit.

% object: theorem-trace-formula-euler-product (blueprint/chap-The-trace-formula/sec-Rationality-and-the-functional-equation/theorem-trace-formula-euler-product/_section.md)
\phantomsection\label{bp92}%
\begin{bpstatementit}{Theorem~12.30}{ (The trace formula as an Euler product)}
Let \(X_0\) be separated and of finite type over \(\mathbf F_q\), let \(\Lambda\in\{\mathbf Q_\ell,E,\overline{\mathbf Q}_\ell\}\), and let \(K\in\Dbc(X,\Lambda)\) carry a Weil structure, with stalk Frobenius maps \(\operatorname{Frob}_{y,m}\) as in Theorem \hyperref[bp73]{12.24}. For each closed point \(x\in|X_0|\), let \(\deg(x)=[\kappa(x):\mathbf F_q]\) be its degree, choose one of the \(\deg(x)\) points \(y\in X_0(\mathbf F_{q^{\deg(x)}})\) lying over \(x\), write \(\bar x=\bar y\) for the corresponding geometric point, and let \(\operatorname{Frob}_x=\operatorname{Frob}_{y,\deg(x)}\) be its action on \(K_{\bar x}\). Another choice of \(y\) replaces \(\operatorname{Frob}_x\) by a conjugate. Define

\[
L(X_0,K,T)
:=\prod_{x\in|X_0|}
\det\bigl(1-T^{\deg(x)}\operatorname{Frob}_x\mid K_{\bar x}\bigr)^{-1}.
\tag{12.6}
\]

For every \(n\), only finitely many factors differ from \(1\) modulo \(T^n\), so the product converges \(T\)-adically in \(1+T\Lambda[[T]]\), and

\[
L(X_0,K,T)
=\prod_i\det\bigl(1-T\operatorname{Frob}_{X,K}\mid H_c^i(X_{\mathrm{pro}\acute{e}t},K)\bigr)^{(-1)^{i+1}}.
\tag{12.7}
\]

In particular, \(L(X_0,K,T)\) is the expansion at zero of a rational function in \(\Lambda(T)\) with value \(1\) at zero. If \(X_0\) is proper, we may replace \(H_c^i\) by \(H^i\).
\end{bpstatementit}

% proof: proof/theorem-trace-formula-euler-product
\begin{proof}[{Proof}]\phantomsection\label{bp93}%
Both sides are of the form \(\exp\) of a power series, and we compare the logarithms coefficient by coefficient. The coefficient of \(T^m\) on the cohomological side is the global trace of \(\operatorname{Frob}^m\) divided by \(m\), and on the side of closed points it is the sum of the local traces over \(X_0(\mathbf F_{q^m})\) divided by \(m\). Hence the theorem is the trace formula for all \(m\) at once.

For convergence, note that \(X_0\) has only finitely many closed points of degree at most any given bound \(n\), because each of them lies under a point of the finite set \(X_0(\mathbf F_{q^{n!}})\). The factor at \(x\) in (12.6) is \(1\) modulo \(T^{\deg(x)}\), so the product converges.

Before comparing logarithms we rewrite the sum over points in Theorem \hyperref[bp73]{12.24} in terms of closed points. A closed point of degree \(d\) has exactly \(d\) points over \(\mathbf F_{q^m}\) when \(d\mid m\), and none otherwise. For instance, \(\operatorname{Spec}\mathbf F_{q^2}\) has one closed point, of degree \(2\), with no \(\mathbf F_q\)-points and two \(\mathbf F_{q^2}\)-points. These \(d\) points form one Frobenius orbit, and the Weil structure gives isomorphisms between the stalks at consecutive points of the orbit. The Frobenius on the stalk at a point of the orbit is then the composite of these isomorphisms taken once around the orbit, starting at that point. Moving the starting point cyclically therefore conjugates this composite by one of these isomorphisms. Hence the operators at the \(\deg(x)\) points are conjugate. In the same way, iterating the Weil structure shows that for \(\deg(x)\mid m\) the operator \(\operatorname{Frob}_{y,m}\) at each of the \(\deg(x)\) points \(y\in X_0(\mathbf F_{q^m})\) over \(x\) is conjugate to \(\operatorname{Frob}_x^{m/\deg(x)}\). Since traces are invariant under conjugation, (12.3) reads

\[
\operatorname{tr}_{\Lambda}(\operatorname{Frob}_{X,K}^m\mid R\Gamma_c(X_{\mathrm{pro}\acute{e}t},K))
=\sum_{\substack{x\in|X_0|\\\deg(x)\mid m}}\deg(x)\,
\operatorname{tr}_{\Lambda}
(\operatorname{Frob}_x^{m/\deg(x)}\mid K_{\bar x}).
\tag{12.8}
\]

We multiply (12.8) by \(T^m/m\) and sum over \(m\geq1\). By Lemma \hyperref[bp89]{12.29}, the left side becomes the logarithm of the right side of (12.7). In the sum over closed points on the right, we substitute \(m=r\deg(x)\). Then the multiplicity \(\deg(x)\) cancels the factor \(\deg(x)\) in \(m\), and Lemma \hyperref[bp89]{12.29} applied to \(\operatorname{Frob}_x\) on \(K_{\bar x}\), with \(T^{\deg(x)}\) in place of \(T\), shows that

\[
\sum_{x\in|X_0|}\sum_{r\geq1}
\operatorname{tr}_{\Lambda}(\operatorname{Frob}_x^r\mid K_{\bar x})
\frac{T^{r\deg(x)}}{r}
=\log L(X_0,K,T).
\]

Each coefficient of \(T\) receives only finitely many contributions, so we may regroup the sums. Exponentiating proves (12.7). Rationality holds because the cohomological product has finitely many factors, each a polynomial or the reciprocal of one, with value \(1\) at zero. If \(X_0\) is proper, compact and ordinary cohomology agree.
\end{proof}

% text: prose-12-30
\phantomsection\label{bp94}%
For the constant sheaf every local factor is \((1-T^{\deg(x)})^{-1}\), and the Euler product becomes the zeta function of \(X_0\). This proves rationality of \(Z(X_0,T)\) over \(\mathbf Q_\ell\). A short argument with linear recurrences descends it to \(\mathbf Q\).

% object: corollary-point-counting-zeta-function (blueprint/chap-The-trace-formula/sec-Rationality-and-the-functional-equation/corollary-point-counting-zeta-function/_section.md)
\phantomsection\label{bp95}%
\begin{bpstatementit}{Corollary~12.31}{ (Point counting and the zeta function)}
Give \(\Lambda_X\) the Weil structure of constant sheaves defined after Definition \bpout{10.18}. For every coefficient ring \(\Lambda\) of Theorem \hyperref[bp73]{12.24} and every \(m\geq1\) we have

\[
\operatorname{tr}_{\Lambda}
(\operatorname{Frob}_{X,\Lambda_X}^m\mid R\Gamma_c(X_{\mathrm{pro}\acute{e}t},\Lambda_X))
=\#X_0(\mathbf F_{q^m})\cdot1_{\Lambda}.
\tag{12.9}
\]

For \(\Lambda=\mathbf Q_\ell,E,\overline{\mathbf Q}_\ell\), this gives

\[
\#X_0(\mathbf F_{q^m})
=\sum_i(-1)^i\operatorname{Tr}
(\operatorname{Frob}_{X,\Lambda_X}^m\mid H_c^i(X_{\mathrm{pro}\acute{e}t},\Lambda_X)),
\]

and

\[
\begin{aligned}
Z(X_0,T)
&:=\exp\left(\sum_{m\geq1}
\#X_0(\mathbf F_{q^m})\frac{T^m}{m}\right)\\
&=\prod_{x\in|X_0|}(1-T^{\deg(x)})^{-1}
=\prod_i\det(1-T\operatorname{Frob}_{X,\Lambda_X}\mid H_c^i(X_{\mathrm{pro}\acute{e}t},\Lambda_X))^{(-1)^{i+1}}.
\end{aligned}
\]

The first two series lie in \(1+T\mathbf Z[[T]]\), and the last equality holds in \(\Lambda[[T]]\). Moreover, \(Z(X_0,T)\) belongs to \(\mathbf Q(T)\).
\end{bpstatementit}

% proof: proof/corollary-point-counting-zeta-function
\begin{proof}[{Proof}]\phantomsection\label{bp96}%
The identities over \(\Lambda\) are the trace formula and the Euler product for the constant sheaf. What needs an argument is that \(Z(X_0,T)\) has rational coefficients, which we deduce at the end from the rationality over \(\mathbf Q_\ell\) by linear algebra.

For \(K=\Lambda_X\), every stalk is \(\Lambda\), and every stalk Frobenius \(\operatorname{Frob}_{y,m}\) is its identity, with categorical trace \(1_\Lambda\). Thus (12.9) follows from Theorem \hyperref[bp73]{12.24}, and the formula over the fields follows from its cohomological form (12.4). By the orbit count in the proof of Theorem \hyperref[bp92]{12.30}, we have \(\#X_0(\mathbf F_{q^m})=\sum_{\deg(x)\mid m}\deg(x)\). Expanding each \((1-T^{\deg(x)})^{-1}\) as a geometric series shows that the product over closed points has integer coefficients. By the orbit count, its logarithm is the series defining \(Z(X_0,T)\). The last equality is Theorem \hyperref[bp92]{12.30} for \(K=\Lambda_X\).

Rationality descends from \(\mathbf Q_\ell\) to \(\mathbf Q\) for the following reason. A power series is rational exactly when its coefficients eventually satisfy a linear recurrence with constant coefficients. The existence of such a recurrence is a system of linear equations, and a system with rational coefficients has a rational solution as soon as it has any solution. More precisely, write \(Z(X_0,T)=\sum_{n\geq0}z_nT^n\), with \(z_n\in\mathbf Z\). Rationality over \(\mathbf Q_\ell\) means that there are integers \(d,N\) and \(b_1,\ldots,b_d\in\mathbf Q_\ell\) such that \(z_n+b_1z_{n-1}+\cdots+b_dz_{n-d}=0\) for all \(n>N\). These are linear equations with rational coefficients in the \(d\) unknowns \(b_i\). Their coefficient rows span a finite-dimensional subspace of \(\mathbf Q^{d+1}\), and hence finitely many of the equations imply all the others. This finite rational linear system has a solution over \(\mathbf Q_\ell\), and hence one over \(\mathbf Q\) by Gaussian elimination. With these rational \(b_i\), the product of \(1+b_1T+\cdots+b_dT^d\in\mathbf Q[T]\) and \(Z(X_0,T)\) is a polynomial. Therefore \(Z(X_0,T)\in\mathbf Q(T)\).
\end{proof}

% text: prose-12-31
\phantomsection\label{bp97}%
For example, let \(X_0=\mathbf A^1_{\mathbf F_q}\). Then \(\#X_0(\mathbf F_{q^m})=q^m\), so that \(Z(X_0,T)=(1-qT)^{-1}\). On the cohomological side, \(H^i_c(X,\mathbf Q_\ell)\) vanishes except for \(i=2\), and \(H^2_c(X,\mathbf Q_\ell)\simeq\mathbf Q_\ell(-1)\) is a line on which geometric Frobenius acts by \(q\). The only factor is therefore \((1-qT)^{-1}\), as it should be.

For \(X_0=\mathbf P^1_{\mathbf F_q}\) the cohomology is \(\mathbf Q_\ell\) in degree \(0\) and \(\mathbf Q_\ell(-1)\) in degree \(2\), so \(Z(\mathbf P^1,T)=1/((1-T)(1-qT))\). This matches \(\#\mathbf P^1(\mathbf F_{q^m})=1+q^m\). For a curve of positive genus the middle cohomology contributes as well, and for the running example a single point count determines it.

% object: example-running-example-trace-formula (blueprint/chap-The-trace-formula/sec-Rationality-and-the-functional-equation/example-running-example-trace-formula/_section.md)
\phantomsection\label{bp98}%
\begin{bpstatement}{Example~12.32}{ (The running example: the trace formula)}
For the elliptic curve \(C_0\colon y^2=x^3-x\) over \(\mathbf F_3\), geometric Frobenius acts on \(H^0(C,\mathbf Q_\ell)=\mathbf Q_\ell\) by \(1\) and on \(H^2(C,\mathbf Q_\ell)=\mathbf Q_\ell(-1)\) by \(3\), and \(H^1(C,\mathbf Q_\ell)\) is a plane (Example \bpout{10.23}). The point-counting formula reads

\[
N_m=1-\operatorname{Tr}\bigl(F^m\mid H^1(C,\mathbf Q_\ell)\bigr)+3^m .
\]

From \(N_1=4\) we get \(\operatorname{Tr}(F\mid H^1)=0\). Duality gives \(\det(F\mid H^1)=3\), so \(\det(1-TF\mid H^1)=1+3T^2\). Hence

\[
Z(C_0,T)=\frac{1+3T^2}{(1-T)(1-3T)}.
\]

The formula predicts \(N_2=1+9+6=16\) and \(N_3=1+27-0=28\), and direct counts over \(\mathbf F_9\) and \(\mathbf F_{27}\) confirm this. In general, for a genus-one curve over \(\mathbf F_q\) the single number \(N_1\) determines \(P_1=1-(1+q-N_1)T+qT^2\), and therefore every \(N_m\).
\end{bpstatement}

% text: prose-12-32
\phantomsection\label{bp99}%
For smooth proper \(X_0\), Poincaré duality relates the factors of the determinant formula in complementary degrees, and this gives a functional equation.

% object: corollary-functional-equation (blueprint/chap-The-trace-formula/sec-Rationality-and-the-functional-equation/corollary-functional-equation/_section.md)
\phantomsection\label{bp100}%
\begin{bpstatementit}{Corollary~12.33}{ (Functional equation)}
Let \(X_0/\mathbf F_q\) be smooth and proper of pure dimension \(d\), and let \(\ell\nmid q\). Put \(H^i=H^i(X_{\mathrm{pro}\acute{e}t},\mathbf Q_\ell)\), let \(F_i\) be the Frobenius operator \(F\) on \(H^i\) (Section \bpout{10.5}), and put

\[
b_i=\dim_{\mathbf Q_\ell}H^i,\qquad \chi=\sum_i(-1)^ib_i,\qquad
D=\prod_i\det(F_i)^{(-1)^i}.
\]

Then \(b_{2d-i}=b_i\), \(d\chi\) is even, \(D^2=q^{d\chi}\), and in \(\mathbf Q_\ell(T)\)

\[
Z\!\left(X_0,\frac1{q^dT}\right)=(-1)^{\chi}D\,T^{\chi}\,Z(X_0,T)
=\varepsilon q^{d\chi/2}T^{\chi}Z(X_0,T),
\qquad \varepsilon=(-1)^\chi Dq^{-d\chi/2}\in\{1,-1\}.
\]
\end{bpstatementit}

% proof: proof/corollary-functional-equation
\begin{proof}[{Proof}]\phantomsection\label{bp101}%
Duality gives the relations \(b_{2d-i}=b_i\) and \(D^2=q^{d\chi}\). If we substitute \(1/(q^dT)\) into each factor of the cohomological expression for \(Z(X_0,T)\), duality turns the factor in degree \(i\) into the factor in degree \(2d-i\) times an explicit scalar.

We apply Theorem \bpout{10.21} with \(K=\mathbf Q_\ell\) and \(X\) proper. It gives \(b_{2d-i}=b_i\), and it shows that the eigenvalues of \(F_{2d-i}\) are the \(q^d/\alpha\) for the eigenvalues \(\alpha\) of \(F_i\), with multiplicities. Hence \(\det F_i\cdot\det F_{2d-i}=q^{db_i}\). Since \((-1)^i=(-1)^{2d-i}\), we get

\[
D^2=\prod_i\bigl(\det F_i\cdot\det F_{2d-i}\bigr)^{(-1)^i}
=q^{d\chi}.
\]

Since \(b_{2d-i}=b_i\), the terms of \(\chi\) for \(i\) and \(2d-i\) are equal and cancel modulo \(2\), so \(\chi\equiv b_d\pmod 2\). If \(d\) is even, \(d\chi\) is even anyway. If \(d\) is odd, the pairing on \(H^d\) of each connected component of \(X\) is perfect and alternating (Corollary \bpout{10.17}), so \(b_d\) is even. In both cases \(d\chi\) is even. Then \(\varepsilon^2=1\) because \(D^2=q^{d\chi}\).

Now consider a single factor. Since \(F_i\) is invertible, its eigenvalues \(\alpha\) are nonzero. Pulling \(-\alpha/(q^dT)\) out of each linear factor gives

\[
\det\bigl(1-q^{-d}T^{-1}F_i\bigr)
=\prod_\alpha\Bigl(-\frac{\alpha}{q^dT}\Bigr)\Bigl(1-\frac{q^d}{\alpha}T\Bigr)
=(-1)^{b_i}\frac{\det F_i}{(q^dT)^{b_i}}\det\bigl(1-TF_{2d-i}\bigr).
\]

By Corollary \hyperref[bp95]{12.31}, and since \(H^i_c=H^i\) as \(X_0\) is proper, \(Z(X_0,T)\) is the product of the \(\det(1-TF_i)^{(-1)^{i+1}}\). We substitute \(1/(q^dT)\) for \(T\) and apply the last display to each factor. Since \((-1)^i=(-1)^{2d-i}\), the factors \(\det(1-TF_{2d-i})^{(-1)^{i+1}}\) multiply to \(Z(X_0,T)\), and the scalar factors multiply to

\[
(-1)^{\chi}D^{-1}(q^dT)^{\chi}=(-1)^{\chi}D\,T^{\chi},
\]

by \(D^2=q^{d\chi}\). The second form of the equation is the definition of \(\varepsilon\).
\end{proof}

% text: prose-12-33
\phantomsection\label{bp102}%
For the elliptic curve of Example \hyperref[bp98]{12.32} we have \(d=1\) and \(b_0=b_2=1\), \(b_1=2\). So \(\chi=0\) and \(D=\det(F_1)^{-1}\det(F_2)=3/3=1\). The corollary says \(Z(C_0,1/(3T))=Z(C_0,T)\), and indeed

\[
\frac{1+3(3T)^{-2}}{(1-(3T)^{-1})(1-T^{-1})}
=\frac{(3T^2+1)/(3T^2)}{(3T-1)(T-1)/(3T^2)}
=\frac{1+3T^2}{(1-3T)(1-T)}.
\]

% object: example-functional-equation-projective-line (blueprint/chap-The-trace-formula/sec-Rationality-and-the-functional-equation/example-functional-equation-projective-line/_section.md)
\phantomsection\label{bp103}%
\begin{bpstatement}{Example~12.34}{ (The functional equation of the projective line)}
For \(X_0=\mathbf P^1_{\mathbf F_q}\) we have \(d=1\), \(b_0=b_2=1\), \(b_1=0\), so \(\chi=2\), and \(D=\det(F_0)\det(F_2)=1\cdot q=q\), so that \(\varepsilon=1\). The corollary predicts \(Z(\mathbf P^1,1/(qT))=qT^2Z(\mathbf P^1,T)\). Directly,

\[
\frac1{(1-(qT)^{-1})(1-T^{-1})}
=\frac{qT^2}{(qT-1)(T-1)}
=\frac{qT^2}{(1-T)(1-qT)}.
\]

Here the factor \(T^{\chi}\) is not trivial, unlike for the elliptic curve, and the constant \(q^{d\chi/2}=q\) comes from the twist \((-1)\) on \(H^2\).
\end{bpstatement}

% object: remark-duality-and-weights (blueprint/chap-The-trace-formula/sec-Rationality-and-the-functional-equation/remark-duality-and-weights/_section.md)
\phantomsection\label{bp104}%
\begin{bpstatement}{Remark~12.35}{ (Duality and weights)}
The functional equation is an identity of rational functions, not of power series at \(T=0\), since \(1/(q^dT)\) is not a power series in \(T\).

What duality does not give is the absolute value of the eigenvalues \(\alpha\). In Chapters \bpout{13}–\bpout{Chapter 19} we prove that the eigenvalues of \(F_i\) have absolute value \(q^{i/2}\) under every embedding of \(\overline{\mathbf Q}_\ell\) into \(\mathbf C\). It follows that the factors \(\det(1-TF_i)\) for different \(i\) have no common zeros, and hence we can read them off from \(Z(X_0,T)\) alone. In Chapter \bpout{19} we use this to show that they are polynomials with integer coefficients that do not depend on \(\ell\).
\end{bpstatement}



\clearpage\phantomsection\hypertarget{bp-outside}{}%
\section*{Outside this export}
\addcontentsline{toc}{section}{Outside this export}
References marked \ensuremath{{}^\ast} are to parts of the blueprint this document leaves out:
\begin{description}
\item[Section 1.3] The function–sheaf dictionary, in 1 Introduction (\bpslug{sec-{}\allowbreak{}The-{}\allowbreak{}function-{}\allowbreak{}sheaf-{}\allowbreak{}dictionary})
\item[Definition 2.61] Derived adic completeness, in Ⅰ The six operations (\bpslug{definition-{}\allowbreak{}derived-{}\allowbreak{}adic-{}\allowbreak{}completeness})
\item[Lemma 3.13] Flat models and generators, in Ⅰ The six operations (\bpslug{lemma-{}\allowbreak{}flat-{}\allowbreak{}models-{}\allowbreak{}generators})
\item[Proposition 3.27] Inverting $\ell$, in Ⅰ The six operations (\bpslug{proposition-{}\allowbreak{}inverting-{}\allowbreak{}ell})
\item[Proposition 5.15] Comparison with the étale derived category, in Ⅰ The six operations (\bpslug{proposition-{}\allowbreak{}comparison-{}\allowbreak{}etale-{}\allowbreak{}derived-{}\allowbreak{}category})
\item[Proposition 7.1] The six operations on $\Dctf$, in Ⅰ The six operations (\bpslug{proposition-{}\allowbreak{}six-{}\allowbreak{}operations-{}\allowbreak{}dctf})
\item[Proposition 7.10] The six operations and change of coefficients, in Ⅰ The six operations (\bpslug{proposition-{}\allowbreak{}six-{}\allowbreak{}operations-{}\allowbreak{}change-{}\allowbreak{}coefficients})
\item[Proposition 7.11] $\Dctf(X,\Qbar)$ as a $2$-colimit, in Ⅰ The six operations (\bpslug{proposition-{}\allowbreak{}dctf-{}\allowbreak{}x-{}\allowbreak{}qbar-{}\allowbreak{}2-{}\allowbreak{}colimit})
\item[Proposition 7.12] Isomorphisms reduce to finite coefficients, in Ⅰ The six operations (\bpslug{proposition-{}\allowbreak{}isomorphisms-{}\allowbreak{}reduce-{}\allowbreak{}finite-{}\allowbreak{}coefficients})
\item[Lemma 10.1] Perfect complexes at the geometric point, in Ⅱ Duality and the trace formula (\bpslug{lemma-{}\allowbreak{}perfect-{}\allowbreak{}complexes-{}\allowbreak{}geometric-{}\allowbreak{}point})
\item[Corollary 10.17] The cup product pairing, in Ⅱ Duality and the trace formula (\bpslug{corollary-{}\allowbreak{}cup-{}\allowbreak{}product-{}\allowbreak{}pairing})
\item[Section 10.5] Weil complexes, in Ⅱ Duality and the trace formula (\bpslug{sec-{}\allowbreak{}Weil-{}\allowbreak{}complexes})
\item[Definition 10.18] Weil structure, in Ⅱ Duality and the trace formula (\bpslug{definition-{}\allowbreak{}weil-{}\allowbreak{}structure})
\item[Definition 10.19] Frobenius operator, in Ⅱ Duality and the trace formula (\bpslug{definition-{}\allowbreak{}frobenius-{}\allowbreak{}operator})
\item[Theorem 10.21] Poincaré duality and Frobenius, in Ⅱ Duality and the trace formula (\bpslug{theorem-{}\allowbreak{}poincare-{}\allowbreak{}duality-{}\allowbreak{}frobenius})
\item[Example 10.23] The running example: duality, in Ⅱ Duality and the trace formula (\bpslug{example-{}\allowbreak{}running-{}\allowbreak{}example-{}\allowbreak{}duality})
\item[Chapter 11] The Lefschetz–Verdier formula, in Ⅱ Duality and the trace formula (\bpslug{chap-{}\allowbreak{}The-{}\allowbreak{}Lefschetz-{}\allowbreak{}Verdier-{}\allowbreak{}formula})
\item[Definition 11.1] Categorical trace, in Ⅱ Duality and the trace formula (\bpslug{definition-{}\allowbreak{}categorical-{}\allowbreak{}trace})
\item[Proposition 11.2] Properties of the trace, in Ⅱ Duality and the trace formula (\bpslug{proposition-{}\allowbreak{}properties-{}\allowbreak{}trace})
\item[Definition 11.5] Correspondences, in Ⅱ Duality and the trace formula (\bpslug{definition-{}\allowbreak{}correspondence})
\item[Definition 11.8] Trace morphism, in Ⅱ Duality and the trace formula (\bpslug{definition-{}\allowbreak{}trace-{}\allowbreak{}morphism})
\item[Section 11.4] Trace classes and local terms, in Ⅱ Duality and the trace formula (\bpslug{sec-{}\allowbreak{}Trace-{}\allowbreak{}classes-{}\allowbreak{}and-{}\allowbreak{}local-{}\allowbreak{}terms})
\item[Definition 11.9] Trace class and local terms, in Ⅱ Duality and the trace formula (\bpslug{definition-{}\allowbreak{}trace-{}\allowbreak{}class-{}\allowbreak{}local-{}\allowbreak{}terms})
\item[Lemma 11.10] Local terms on the point, in Ⅱ Duality and the trace formula (\bpslug{lemma-{}\allowbreak{}local-{}\allowbreak{}terms-{}\allowbreak{}point})
\item[Definition 11.11] Action of a correspondence on cohomology, in Ⅱ Duality and the trace formula (\bpslug{definition-{}\allowbreak{}action-{}\allowbreak{}correspondence})
\item[Theorem 11.13] Proper pushforward and Lefschetz–Verdier, in Ⅱ Duality and the trace formula (\bpslug{theorem-{}\allowbreak{}proper-{}\allowbreak{}pushforward-{}\allowbreak{}lefschetz-{}\allowbreak{}verdier})
\item[Proposition 11.19] Additivity of trace classes, in Ⅱ Duality and the trace formula (\bpslug{proposition-{}\allowbreak{}additivity-{}\allowbreak{}trace-{}\allowbreak{}classes})
\item[Chapter 13] Weights, in Ⅲ Weights and the Riemann hypothesis (\bpslug{chap-{}\allowbreak{}Weights})
\item[Chapter 16] The Fourier transform, in Ⅲ Weights and the Riemann hypothesis (\bpslug{chap-{}\allowbreak{}The-{}\allowbreak{}Fourier-{}\allowbreak{}transform})
\item[Lemma 16.1] The Artin–Schreier character sheaf, in Ⅲ Weights and the Riemann hypothesis (\bpslug{lemma-{}\allowbreak{}artin-{}\allowbreak{}schreier-{}\allowbreak{}character-{}\allowbreak{}sheaf})
\item[Chapter 19] Proof of the Weil conjectures, in Ⅲ Weights and the Riemann hypothesis (\bpslug{chap-{}\allowbreak{}Proof-{}\allowbreak{}of-{}\allowbreak{}the-{}\allowbreak{}Weil-{}\allowbreak{}conjectures})
\item[Appendix D] Cohomology of torsion sheaves on curves, in Appendices (\bpslug{chap-{}\allowbreak{}Cohomology-{}\allowbreak{}of-{}\allowbreak{}torsion-{}\allowbreak{}sheaves-{}\allowbreak{}on-{}\allowbreak{}curves})
\item[Proposition D.4] Cohomology of a smooth affine curve, in Appendices (\bpslug{proposition-{}\allowbreak{}cohomology-{}\allowbreak{}smooth-{}\allowbreak{}affine-{}\allowbreak{}curve})
\end{description}
\clearpage\phantomsection%
\section*{References}
\addcontentsline{toc}{section}{References}
\bpbibitem{BBD82}{bp107}{A. Beilinson, J. Bernstein, and P. Deligne. \emph{Faisceaux pervers}, volume 100 of \emph{Astérisque}. 1982.}

\bpbibitem{BS15}{bp108}{B. Bhatt and P. Scholze. The pro-étale topology for schemes. \emph{Astérisque}, 369:99–201, 2015. \href{https://people.mpim-bonn.mpg.de/scholze/proetale.pdf}{Author's manuscript}.}

\bpbibitem{Con07}{bp109}{B. Conrad. Deligne's notes on Nagata compactifications. \emph{Journal of the Ramanujan Mathematical Society}, 22:205–257, 2007. \href{https://math.stanford.edu/\~conrad/papers/nagatafinal.pdf}{Author's manuscript}.}

\bpbibitem{Del74}{bp110}{P. Deligne. La conjecture de Weil. I. \emph{Publ. Math. IHÉS}, 43:273–307, 1974.}

\bpbibitem{Del80}{bp111}{P. Deligne. La conjecture de Weil. II. \emph{Publ. Math. IHÉS}, 52:137–252, 1980.}

\bpbibitem{Den22}{bp112}{C. Deninger. There is no “Weil-”cohomology theory with real coefficients for arithmetic curves. \emph{Ann. Sc. Norm. Super. Pisa Cl. Sci.}, 2022. To appear. arXiv:\href{https://arxiv.org/abs/2204.02714}{2204.02714}.}

\bpbibitem{Deu41}{bp113}{M. Deuring. Die Typen der Multiplikatorenringe elliptischer Funktionenkörper. \emph{Abh. Math. Sem. Univ. Hamburg}, 14:197–272, 1941.}

\bpbibitem{Eke90}{bp114}{T. Ekedahl. On the adic formalism. In \emph{The Grothendieck Festschrift, Vol. II}, volume 87 of \emph{Progress in Mathematics}, pages 197–218. Birkhäuser, 1990.}

\bpbibitem{Has36}{bp115}{H. Hasse. Zur Theorie der abstrakten elliptischen Funktionenkörper. I, II, III. \emph{J. Reine Angew. Math.}, 175:55–62, 69–88, 193–208, 1936.}

\bpbibitem{Hum72}{bp116}{J. E. Humphreys. \emph{Introduction to Lie Algebras and Representation Theory}, volume 9 of \emph{Graduate Texts in Mathematics}. Springer, 1972.}

\bpbibitem{Ill94}{bp117}{L. Illusie. Autour du théorème de monodromie locale. \emph{Astérisque}, 223:9–57, 1994. §4. \href{https://www.numdam.org/item/AST_1994__223__9_0/}{Original article}.}

\bpbibitem{Kat88}{bp118}{N. M. Katz. Travaux de Laumon. \emph{Astérisque}, 161–162:105–132, 1988. Séminaire Bourbaki, exp. 691.}

\bpbibitem{KW01}{bp119}{R. Kiehl and R. Weissauer. \emph{Weil conjectures, perverse sheaves and \(\ell\)-adic Fourier transform}, volume 42 of \emph{Ergebnisse der Mathematik}. Springer, 2001.}

\bpbibitem{Lau87}{bp120}{G. Laumon. Transformation de Fourier, constantes d'équations fonctionnelles et conjecture de Weil. \emph{Publ. Math. IHÉS}, 65:131–210, 1987.}

\bpbibitem{Mil86a}{bp121}{J. S. Milne. \href{https://www.jmilne.org/math/xnotes/AVs.pdf}{Abelian Varieties}. In \emph{Arithmetic Geometry}. 1986. Corrected online version.}

\bpbibitem{Mil86b}{bp122}{J. S. Milne. \href{https://www.jmilne.org/math/xnotes/JVs.pdf}{Jacobian Varieties}. In \emph{Arithmetic Geometry}. 1986. Corrected online version.}

\bpbibitem{Ran39}{bp123}{R. A. Rankin. Contributions to the theory of Ramanujan's function \(\tau(n)\) and similar arithmetical functions. II. \emph{Proc. Cambridge Philos. Soc.}, 35:357–372, 1939.}

\bpbibitem{SGA4}{bp124}{M. Artin, A. Grothendieck, and J.-L. Verdier, editors. \emph{Théorie des topos et cohomologie étale des schémas}, volumes 269, 270, 305 of \emph{Lecture Notes in Mathematics}. Springer, 1972–1973. SGA 4.}

\bpbibitem{SGA4½}{bp106}{P. Deligne. \emph{Cohomologie étale}, volume 569 of \emph{Lecture Notes in Mathematics}. Springer, 1977. SGA 4½; \emph{Théorèmes de finitude}, §3, and \emph{Rapport sur la formule des traces}, §3, and \emph{Sommes trigonométriques}, §1.}

\bpbibitem{SGA5}{bp105}{A. Grothendieck, editor. \emph{Cohomologie \(\ell\)-adique et fonctions \(L\)}, volume 589 of \emph{Lecture Notes in Mathematics}. Springer, 1977. SGA 5.}

\bpbibitem{SGA7II}{bp125}{P. Deligne and N. Katz, editors. \emph{Groupes de monodromie en géométrie algébrique}, volume 340 of \emph{Lecture Notes in Mathematics}. Springer, 1973. SGA 7 II; Exposé XIII, (2.1.7).}

\bpbibitem{Stacks}{bp126}{The Stacks Project Authors. \href{https://stacks.math.columbia.edu}{The Stacks Project}. Individual results are cited by their stable tags in the text.}

\bpbibitem{Var05}{bp127}{Y. Varshavsky. Lefschetz–Verdier trace formula and a generalization of a theorem of Fujiwara. 2005. Version 2. arXiv:\href{https://arxiv.org/abs/math/0505564}{math/0505564}. \href{https://arxiv.org/abs/math/0505564v2}{Author's preprint}.}

\bpbibitem{Wei48a}{bp128}{A. Weil. \emph{Sur les courbes algébriques et les variétés qui s'en déduisent}. Hermann, Paris, 1948.}

\bpbibitem{Wei48b}{bp129}{A. Weil. \emph{Variétés abéliennes et courbes algébriques}. Hermann, Paris, 1948.}

\bpbibitem{Wei49}{bp130}{A. Weil. Numbers of solutions of equations in finite fields. \emph{Bull. Amer. Math. Soc.}, 55:497–508, 1949.}

\end{document}
